\documentclass[11pt]{article}
\usepackage[T1]{fontenc}
\usepackage{lmodern,amsmath,amssymb,amsthm,mathtools,booktabs,longtable,array}
\usepackage[margin=1.05in]{geometry}
\usepackage{microtype}
\usepackage[colorlinks=true,linkcolor=blue!50!black,urlcolor=blue!60!black,citecolor=blue!50!black]{hyperref}
\usepackage{xcolor}
\hypersetup{pdftitle={A determinant argument for four odd zeta values},pdfsubject={Completed reduction, finite certificates, and historical context}}
\usepackage{enumitem}
\setlist{itemsep=3pt,topsep=5pt}
\newtheorem{theorem}{Theorem}[section]
\newtheorem{lemma}[theorem]{Lemma}
\newtheorem{proposition}[theorem]{Proposition}
\newtheorem{corollary}[theorem]{Corollary}
\theoremstyle{definition}\newtheorem{definition}[theorem]{Definition}
\theoremstyle{remark}\newtheorem{remark}[theorem]{Remark}
\newcommand{\Q}{\mathbb Q}\newcommand{\R}{\mathbb R}\newcommand{\C}{\mathbb C}\newcommand{\Z}{\mathbb Z}\newcommand{\N}{\mathbb N}
\newcommand{\Res}{\operatorname{Res}}\newcommand{\rank}{\operatorname{rank}}\newcommand{\lc}{\operatorname{lc}}
\newcommand{\ord}{\operatorname{ord}}\newcommand{\dd}{\,\mathrm d}\newcommand{\eps}{\varepsilon}
\newcommand{\floor}[1]{\left\lfloor#1\right\rfloor}
\newcommand{\abs}[1]{\left|#1\right|}
\newcommand{\norm}[1]{\left\|#1\right\|}
\newcommand{\supp}{\operatorname{supp}}
\DeclareMathOperator*{\liminfN}{lim\,inf}
\title{A determinant argument for four odd zeta values\\[3pt]\large Completed reduction and finite certificates}
\author{}
\date{Revised 15 October 2026}
\begin{document}
\maketitle
\begin{abstract}
We develop a determinant argument concerning the four values
$\zeta(7)$, $\zeta(9)$, $\zeta(11)$ and $\zeta(13)$, proving that they
cannot all be rational.  A rational-function moment family
has a fourth Taylor coefficient involving exactly these four values.
A finite Pearson-functional classification gives uniform nonvanishing of
an almost full Hankel minor, including after arbitrary complex substitution
for the four zeta values.  A tilted contour and an explicitly specified
bounded-root polynomial basis give an upper exponent for every such minor.
Local $p$-adic jet estimates give a lower exponent under simultaneous
rationality.  The numerical parts are reduced to finite, reproducible
integer/rational and interval-arithmetic certificates, with normalized
gap $0.03706678$. The mathematical reduction to finite certificates has
been checked in Lean; eight isolated numerical admissions are independently
verified by exact rational and rigorous ball-arithmetic computations.
\end{abstract}
\noindent\textbf{Verification status.}
This revision incorporates the completed formal reduction and reproducibility
audit. The supplied-certificate theorem retains eight explicitly identified
finite-calculation admissions; it is not an admission-free formalization.
The manuscript has not undergone independent peer review.
\clearpage
\begingroup
\setcounter{tocdepth}{1}
\tableofcontents
\endgroup
\clearpage

\section{Introduction}\subsection{Historical background}
Euler's evaluation of the even zeta values,
\[
 \zeta(2k)=(-1)^{k+1}\frac{B_{2k}(2\pi)^{2k}}{2(2k)!},
\]
places their arithmetic in a different setting from that of the odd values;
see \cite[\S25.6]{DLMF}. Ap\'ery's 1978 discovery, published in
\cite{Apery}, established the irrationality of $\zeta(3)$ through specially
structured rational approximations. Beukers subsequently gave an integral
proof \cite{Beukers}. These arguments illustrate a recurrent principle:
one must construct an expression that is both analytically small and
arithmetically constrained, while also proving that it does not vanish.
The last condition is indispensable, even when strong denominator bounds
and decay estimates are available.

Rivoal \cite{Rivoal00}, and Ball and Rivoal \cite{BR}, established that
infinitely many odd zeta values are irrational. More quantitatively, their
work gives a logarithmically growing lower bound for the dimension over
$\Q$ of the span of $1,\zeta(3),\zeta(5),\ldots,\zeta(s)$ as the odd
integer $s$ tends to infinity. This is a linear-independence statement,
stronger than a corresponding count of irrational entries. Fischler's
Bourbaki account \cite{Fischler04} explains the connections among rational
approximation, hypergeometric series, and the irrationality criteria used
in this development.

Explicit short lists pose a related but distinct problem. Rivoal's
nine-value result \cite{Rivoal02}, announced in 2001, concerns
$\zeta(5),\zeta(7),\ldots,\zeta(21)$. Zudilin proved in 2001 that at least
one of $\zeta(5),\zeta(7),\zeta(9),\zeta(11)$ is irrational
\cite{Zudilin01}. His later systematic treatment \cite{Zudilin04} develops
hypergeometric constructions and their arithmetic, including the removal
of common prime factors. Integral and series representations of such
linear forms are also developed by Nesterenko \cite{Nesterenko03}.

Subsequent work has pursued both more elementary proofs and stronger
asymptotic counts. Zudilin \cite{Zudilin18} obtained an elementary proof
that one of the odd values from $\zeta(5)$ through $\zeta(25)$ is irrational,
using the prime number theorem and Stirling's formula among its analytic
ingredients. Fischler, Sprang and Zudilin \cite{FSZ} proved a lower bound
$2^{(1-\varepsilon)\log s/\log\log s}$ for the number of irrational odd
zeta values up to $s$, for sufficiently large $s$. Lai and Yu \cite{LY}
improved this to $(1.192507\ldots-\varepsilon)\sqrt{s/\log s}$.
Lai's subsequent refinement \cite{Lai26} gives
$1.284\sqrt{s/\log s}$ among $\zeta(3),\zeta(5),\ldots,\zeta(s-1)$ for
sufficiently large even $s$, combining the elimination method with
Zudilin's prime-divisor improvement. These are counts of irrational
values; they should not be confused with lower bounds of the same size
for the dimension of their rational span.

The fixed set studied here is $\{\zeta(7),\zeta(9),\zeta(11),\zeta(13)\}$.
It excludes $\zeta(5)$ and includes $\zeta(13)$, so the result of
\cite{Zudilin01} alone does not imply the desired conclusion. Nor does an
asymptotic counting theorem automatically settle this particular short
interval. The comparison concerns the supports of the assertions, not a
claim to subsume those earlier theorems. The discussion above records relevant precedents rather than an
exhaustive priority survey.

\subsection{The determinant construction and its verification}
The starting point is a factorial-normalized rational function, in the
tradition of \cite{BR,Zudilin04}. Its fourth Taylor coefficient shifts the
surviving odd pole orders $3,5,7,9$ to the arguments $7,9,11,13$.
The present argument uses a growing moment matrix rather than requiring
one prescribed scalar linear form to be nonzero. A finite Pearson
classification and a truncated Hankel extension force its rank to be
within four of its size, uniformly under arbitrary complex substitutions
for the four target zeta values. One then selects an almost full nonzero
minor. Uniform estimates for every minor of that size ensure that the
selection creates no unproved size condition.

On the analytic side, a tilted contour and an explicit bounded-root basis
give a uniform upper exponent. On the arithmetic side, local square-minor
valuations are combined by a finite dual inequality; rational boxes control
all floor values on whole parameter cells. The use of $p$-adic valuations
here concerns rational numbers and the product formula. The theorem is
about the classical real zeta function, not a $p$-adic zeta function.

This version incorporates the completed Round~29 reduction and final audit
\cite{Project}. Lean~4 \cite{Lean4}, Mathlib \cite{Mathlib}, and a pinned
prime-number-theorem development \cite{PNT} check the mathematical
implications from explicit finite certificates to the classical zeta
conclusion. Eight isolated admissions supply the finite numerical families;
separate programs check their data using exact rational arithmetic and Arb
ball arithmetic \cite{Arb}. Thus the result has a completed formal reduction
with externally checked finite inputs. It is not an admission-free Lean
formalization. Section~\ref{sec:verification} and
Appendix~\ref{app:reproduction} specify the boundary and reproduction steps.

\subsection{Statement and rates}

For an integer $s>1$, write $\zeta(s)=\sum_{m=1}^\infty m^{-s}$.
With the finite certificates specified below, the conclusion is the following.
\begin{theorem}\label{thm:main}
At least one of $\zeta(7),\zeta(9),\zeta(11),\zeta(13)$ is irrational.
\end{theorem}
The argument is by contradiction.  With $M=100n$, we construct a nonzero
rational number $D_n$ under the hypothesis that all four values are rational.
Its definition involves a minor selected from a matrix of dimension
$N=\floor{361n/2}+1$; the selected row and column sets need not be the same.
Uniformly in these sets, we obtain
\begin{align}
\limsup_{n\to\infty}M^{-2}\log\abs{D_n}
 &\le A_*=-\frac{577853339}{50000000},\label{eq:main-upper}\\
\liminf_{n\to\infty}M^{-2}\log\abs{D_n}
 &\ge P_*=-\frac{288}{25}.\label{eq:main-lower}
\end{align}
The arithmetic assertion is stronger than the displayed liminf statement:
its lower bound holds for every sufficiently large $n$. The exact difference is
\begin{equation}\label{eq:gap}
 P_*-A_* = \frac{1853339}{50000000}
 =0.03706678>0.
\end{equation}
The nonvanishing argument is algebraic and uniform in the putative rational
values.  It avoids any assertion that a particular principal determinant is
nonzero.  The estimates use every almost full minor, so the choice permitted
by the rank argument creates no additional size problem.

Sections~\ref{sec:family}--\ref{sec:rank} give the exact family and the
nonvanishing argument.  Sections~\ref{sec:real}--\ref{sec:real-cert} prove
and certify the real estimate.  Sections~\ref{sec:local}--\ref{sec:prime}
give the arithmetic estimate and its certificate. The subsequent sections apply
the rational product formula and describe verification. All logarithms are natural.  Implied constants
may depend on the fixed endpoint list and the fixed certificate, but never
on $n$, on a selected minor, or on a prime in the indicated range.

\section{The rational-function family}\label{sec:family}
The construction uses the symmetry and partial-fraction mechanism familiar
from hypergeometric approaches to odd zeta values \cite{Zudilin04}.
The endpoint tuple and the determinant estimates below are the specific
inputs of the present argument.
\subsection{Endpoints and normalization}
Put
\begin{equation}\label{eq:endpoints}
 T=(63,60,58,57,54,52,49,47,44,42),\qquad
 (a_1,\ldots,a_{15})=n(67,67,67,67,67,T).
\end{equation}
Set $b=100n$, $a_0=67n$, and $\beta=\sum_{i=1}^{15}a_i-5b=361n$.
For any integer $a\ge0$, define
\[
 D_a(x)=\prod_{j=-a}^{a}(x+j),\qquad
 \mathcal N_n=\frac{\prod_{t\in T}(2tn)!}{((33n)!)^{10}}.
\]
Our weight is
\begin{equation}\label{eq:weight}
 w_n(x)=\mathcal N_n\frac{2xD_b(x)^5}{\prod_{i=1}^{15}D_{a_i}(x)}.
\end{equation}
It is odd, has degree $-2\beta-9=-722n-9$, and has no poles outside
$[-63n,63n]$.  A nonzero integer $j$ in this interval has pole order
\[
 o_j=\#\{t\in T:|j|\le tn\}\le10.
\]
At the origin the order is nine, due to the explicit factor $x$ in the
numerator.  The integers $a_0+1,\ldots,b$ are zeros of order five.
These statements follow by counting the factors $x-j$ in numerator and
denominator.  In particular, the five copies of $D_{67n}$ cancel the
numerator zeros throughout the range $|j|\le67n$ before the remaining
ten denominator factors are counted.

\subsection{Moment functionals and the four zeta values}
If $\deg P\le\beta+3$ and $0\le h\le4$, define
\begin{equation}\label{eq:F}
 F_h(P)=\sum_{m=b+1}^{\infty}[u^h]\{w_n(m+u)P((m+u)^2)\}.
\end{equation}
At the maximum permitted degree the rational function inside the braces
has degree at most $-3$.  Termwise Taylor differentiation up to order four
is therefore justified by absolute, locally uniform convergence.  The
fivefold zeros allow the lower summation limit to be replaced by $a_0+1$.
Write the partial fractions as
\begin{equation}\label{eq:pf}
 w_n(x)P(x^2)=\sum_j\sum_{k=1}^{o_j}\frac{c_{k,j}(P)}{(x-j)^k},
 \qquad B_k(P)=\sum_jc_{k,j}(P),
\end{equation}
where the center order is interpreted as nine and missing coefficients
are zero.  Reflection gives
$c_{k,-j}=(-1)^{k+1}c_{k,j}$, so $B_k=0$ for even $k$.
The residue at infinity vanishes, and hence $B_1=0$ as well.
For $H_m^{(s)}=\sum_{r=1}^m r^{-s}$, with $H_0^{(s)}=0$, set
\begin{equation}\label{eq:C}
 C_h(P)=-(-1)^h\sum_{j,k}c_{k,j}(P)
       \binom{k+h-1}{h}H_{a_0-j}^{(k+h)}.
\end{equation}
The apparent $\zeta(1)$ term when $h=0$ cancels because $B_1=0$;
the finite harmonic formula remains meaningful.  Summing the other
partial fractions gives
\begin{equation}\label{eq:zeta-form}
 F_4(P)=C_4(P)+\sum_{s\in\{7,9,11,13\}}
                  \binom{s-1}{4}B_{s-4}(P)\zeta(s).
\end{equation}
Thus exactly the four target odd zeta values occur.  Each $B_k$ and $C_h$
takes rational values on $\Q[Y]$.

Under the contrary rationality hypothesis choose $Q\in\N_{>0}$ with
$Q\zeta(s)\in\Z$ for the four target exponents.  Put $\gamma_s=\zeta(s)$
at those exponents and $\gamma_s=0$ for the other integers $5\le s\le14$.
Then
\begin{equation}\label{eq:arithmetic-form}
 F_4(P)=\sum_{j,k}c_{k,j}(P)\binom{k+3}{4}
              \left(\gamma_{k+4}-H_{a_0-j}^{(k+4)}\right).
\end{equation}
All harmonic indices lie in $[0,130n]$.  Equation~\eqref{eq:arithmetic-form}
is an identity in the proper degree range; the terms assigned zero
$\gamma_s$ have zero total coefficient.  This representation is especially
useful because its denominator losses can be localized at the poles.

\section{Pearson relations and the initial moment space}\label{sec:pearson}
\subsection{A difference operator}
Use $z=2x$ and $X=z^2=4Y$.  The rescaling is a diagonal change of polynomial
basis and does not affect any rank. It is used only for the Pearson
argument. All analytic and arithmetic determinants later in the paper use
$Y=x^2$ and $E_i(Y)$; evaluating $E_i(4x^2)$ would change their normalization.
The formal corrected-coordinate adapter proves this identification for
every selected row and column set.  Define
\begin{align}
 A(z)&=(z+2b+2)^5\prod_{i=1}^{15}(z-2a_i),\label{eq:A}\\
 \mathcal Dq(z^2)&=
 \frac{A(z)q((z+1)^2)-A(-z)q((z-1)^2)}{z}.\label{eq:D}
\end{align}
The numerator in~\eqref{eq:D} is odd, so the quotient is an even polynomial.
The polynomial $A$ is monic of degree twenty, and its next coefficient is
$10-2\beta$.  Comparing the two highest potentially surviving powers gives
\begin{equation}\label{eq:D-leading}
 \deg\mathcal D(X^k)=k+9,\qquad
 \lc(\mathcal D(X^k))=4(k-\beta+5)
 \quad(0\le k\le\beta-6).
\end{equation}
The restrictions on $k$ ensure that the displayed coefficient is nonzero.

\begin{lemma}\label{lem:pearson}
The nine functionals
\begin{equation}\label{eq:nine}
 B_3,B_5,B_7,B_9,C_0,C_1,C_2,C_3,C_4
\end{equation}
annihilate $\mathcal D(X^k)$ for $0\le k\le\beta-6$, after the
change $X=4Y$.
\end{lemma}
\begin{proof}
Let $g(x)=w_n(x)/x$.  Direct cancellation of consecutive products gives
\[
 \frac{g(x+1)}{g(x)}=\frac{A(2x)}{A(-2x-2)}.
\]
Consequently, with $H(x)=g(x)A(2x)q((2x+1)^2)$,
\[
 w_n(x)\mathcal Dq(4x^2)=\tfrac12\bigl(H(x)-H(x-1)\bigr).
\]
For $\deg q=k\le\beta-6$, the degree of $H$ is at most $-2$.
Summation over $x=b+1+u,b+2+u,\ldots$ telescopes to $-H(b+u)/2$.
The latter has a zero of order at least five at $u=0$.  It follows that
the five $F_h$ vanish on the relation.  A rational difference
$H(x)-H(x-1)$ has zero sum of Laurent coefficients at every fixed pole
order: translation only permutes the poles.  Thus every $B_k$ vanishes
on the relation.  The partial-fraction identities for $F_h$ then show
that the five $C_h$ vanish as well.
\end{proof}

\subsection{Independence by an endpoint recurrence}
We prove independence in a slightly larger parameter family, solely to
make a convenient seed available.  Temporarily let $0\le a_i\le b$ be
arbitrary integers, with $\beta=\sum a_i-5b\ge6$ and $b>5$.
The scalar normalization is irrelevant to independence and may be omitted.
The definition of $C_h$ can use the common summation boundary $b+1$;
for the target tuple it agrees with~\eqref{eq:C}, since the intervening
terms have fivefold zeros.

\begin{lemma}\label{lem:independent}
For every such tuple, the nine functionals in~\eqref{eq:nine} are
linearly independent on polynomials of degree at most eight.
\end{lemma}
\begin{proof}
If a positive endpoint $a_i$ is decreased by one, the weight is multiplied,
up to a nonzero scalar, by $Y-a_i^2$.  Therefore the first nine moments
of each of the functionals are transformed by the same companion matrix
for multiplication by $Y-a_i^2$.  To compute its determinant, use the
degree-nine relation $\mathcal D1$ to express the ninth moment in terms
of the first nine.  The characteristic polynomial of the companion
matrix is proportional to $\mathcal D1$, in the $X$ coordinate.
It is nonzero at $X=4a_i^2$, because
\[
 (\mathcal D1)(4a_i^2)=-\frac{A(-2a_i)}{2a_i}\ne0.
\]
Indeed $A(2a_i)=0$, while $-2a_i$ equals neither $-2b-2$ nor any
nonnegative root of $A$.  The leading coefficient $4(5-\beta)$ is
nonzero.  Thus the transformation is invertible whenever the old and
new tuples remain in the stated range.

Every admissible tuple can be connected to the all-$b$ tuple by increasing
endpoints one at a time, never decreasing $\beta$.  It suffices to check
the seed consisting of six endpoints equal to $b$ and nine equal to zero;
it has $\beta=b$.  Its weight is a nonzero constant times
$2/(x^8D_b(x))$.  On the monomials $1,Y,\ldots,Y^8$, the four $B$ columns
have a nonsingular anti-triangular block in the first four rows: their
successive highest center poles have orders $9,7,5,3$.  These columns
vanish in the last five rows.

For the last five rows the rational functions are
$2x^{2\ell}/D_b(x)$, $0\le\ell\le4$.  Let
$R(x)=(x-b)_{2b}=\prod_{j=-b}^{b-1}(x+j)$.
There is a unique polynomial $Q_\ell$ in $x(x-1)$ of degree $\ell$ such that
\begin{equation}\label{eq:seed-antidifference}
 (x+b)Q_\ell(x)-(x-b)Q_\ell(x+1)=2x^{2\ell}.
\end{equation}
To verify existence, apply the left-hand operator to $[x(x-1)]^h$.
Its output is even of degree $2h$, with leading coefficient $2(b-h)$.
The resulting triangular system is invertible for $0\le h\le4$.
In particular $\lc Q_\ell=1/(b-\ell)$.  Division by $D_b(x)$ now shows
that the sum from $b+1+u$ to infinity equals
$Q_\ell(b+1+u)/R(b+1+u)$.  The boundary term at infinity vanishes.
The first five Taylor coefficients of these five functions form an
invertible matrix: multiplication by $1/R(b+1+u)$ is invertible on
five-jets, as is substitution of $x(x-1)$ at $x=b+1$, whose derivative
is $2b+1\ne0$.  The polynomials $Q_0,\ldots,Q_4$ form a basis by their
nonzero leading coefficients.  This proves the remaining five-by-five
block invertible, and hence proves independence at the seed and at every
admissible tuple.
\end{proof}

\section{Finite Pearson functionals and uniform nonvanishing}\label{sec:rank}
\subsection{Classification of finite solutions}
A linear functional on $\C[z]$ is called finite if it factors through
$\C[z]/(p)$ for some nonzero polynomial $p$.  Equivalently, it is a finite
sum of evaluation derivatives at finitely many points.  Its local component
at a point is a jet; multiplication by a polynomial nonzero at that point
is invertible on the local jet space.  Multiplication by a zero of order
$d$ lowers a nonzero highest derivative order $h$ to $h-d$, or kills
the jet if $h<d$.

\begin{proposition}\label{prop:finite-classification}
For the endpoints~\eqref{eq:endpoints}, the space of finite functionals $U$
satisfying
\begin{equation}\label{eq:finite-pearson}
 U(f(-z))=-U(f(z)),\qquad
 U\bigl(A(z)q((z+1)^2)\bigr)=0\quad(q\in\C[X])
\end{equation}
has dimension seven.  Each solution is supported on the even integers
between $-134n$ and $134n$ and is determined by its odd center jets of
orders $1,3,\ldots,13$.
\end{proposition}
\begin{proof}
The polynomials $q((z+1)^2)$ are precisely the polynomials invariant under
$\tau(z)=-2-z$.  Consequently the second identity says that $AU$ is odd
under $\tau$.  The first says that $U$ is odd under $\sigma(z)=-z$.
Away from a zero of $A$, these identities propagate a support point by
$\sigma\tau(z)=z+2$.  A finite nonempty chain must start at $-a$ and end
at $a'$ for roots $a,a'$ of $A$, with $a+a'\ge0$.

The roots of $A$ are $-200n-2$ and the positive even integers $2a_i$.
A pair involving the negative root has negative sum, since $2a_i\le134n$.
Thus both endpoints arise from positive roots.  Every chain lies on the
even integer lattice, in the indicated interval, and crosses zero.
There can be no finite cycle under translation by two over $\C$.

On the nonnegative half-lattice, the local equation at consecutive points
can be written as equality of the translated jets obtained by multiplying
at $2j$ by $A(z)$ and at $2j+2$ by $A(-z)$.  The latter factor is a unit
because $A(-2j-2)\ne0$ for $0\le j\le67n$.  Thus a center jet determines
every later jet uniquely.  Its highest order decreases by the multiplicity
of each positive root crossed.  There are fifteen such roots counted with
multiplicity.  A highest center order below fifteen is killed by the
terminal endpoint; an order at least fifteen survives and cannot yield
finite support.  Oddness allows only odd orders, giving exactly
$1,3,\ldots,13$.  Conversely, each such center jet propagates and
terminates; reflection provides the negative half and satisfies the same
local equations.  This proves both the dimension and the claimed
parametrization.  In particular, the five coincident roots at $134n$
lower the order by five; distinctness is not assumed.
\end{proof}

\begin{proposition}\label{prop:finite-restrictions}
For $\ell_U(P)=U(zP(z^2))$, the restrictions of all solutions in
Proposition~\ref{prop:finite-classification} to polynomials of degree at
most eight are exactly
\begin{equation}\label{eq:finite-span}
 \operatorname{span}_{\C}\{B_3,B_5,B_7,B_9,C_1,C_3\}.
\end{equation}
Here and below the fixed change $X=4Y$ is understood.  This subspace has
dimension six and contains no functional with nonzero $C_4$ coordinate
in the basis~\eqref{eq:nine}.
\end{proposition}
\begin{proof}
The even rational function $g(z)=w_n(z/2)/z$ satisfies
\[
 g(z+2)/g(z)=A(z)/A(-z-2)
\]
and has a center pole of order ten.  At every even integer use the same
even Laurent germ in the local coordinate.  Internal residues of $g$
times any such germ satisfy the local Pearson propagation equation.
The germs $1,u^2,u^4,u^6,u^8$ give, after multiplication of the test
polynomial by $z$, nonzero scalar multiples of $B_1,B_3,B_5,B_7,B_9$.
Their highest center derivative orders are $9,7,5,3,1$.
Nonnegative-power germs are merely local data; a globally defined
periodic function with these exact germs is not required.

Two further choices are the globally periodic kernels
$(\pi/2)^2\csc^2(\pi z/2)$ and $(\pi/2)^4\csc^4(\pi z/2)$.
Their highest center derivative orders are eleven and thirteen.  Hence
the seven residue functionals just described form a basis of the finite
solution space, by their distinct highest center orders.

It remains to identify the two additional restrictions.  This is most
transparent in the physical $x$ coordinate.  Put
\[
 K_2(x)=\pi^2\csc^2(\pi x),\qquad
 K_4(x)=\pi^4\csc^4(\pi x).
\]
Their principal parts at an integer are respectively $u^{-2}$ and
$u^{-4}+(2\pi^2/3)u^{-2}$.  Integrate
$w_n(x)P(x^2)K_h(x)$ on expanding rectangles with vertical sides
at half-integers.  For $\deg P\le8$, rational decay makes the vertical
integrals tend to zero, while the horizontal integrals decay
exponentially.  The sum of all residues is therefore zero.  The zeros
of order five at $a_0<|m|\le b$ remove every intervening residue.
At the remaining external integers reflection doubles the positive
half of the sum.  Thus the two internal residue functionals, in this
coordinate, are
\begin{equation}\label{eq:kernel-residues}
 -2F_1(P),\qquad -2F_3(P)-\frac{4\pi^2}{3}F_1(P).
\end{equation}
Internal residues must include every integer $|m|\le a_0$, even when
$w_n$ is regular there: the trigonometric kernel still has a pole.
In particular the nodes $63n<|m|\le67n$ cannot be omitted.

The partial-fraction formula expresses $F_1$ and $F_3$ as $C_1$ and
$C_3$, respectively, plus linear combinations of the four nonzero $B$
directions.  Only even zeta values occur in these combinations, and
their values are irrelevant to the span.  Also $B_1=0$ in the proper
range.  Equation~\eqref{eq:kernel-residues} is triangular with nonzero
coefficients on $C_1,C_3$.  Lemma~\ref{lem:independent} now proves the
dimension and the final exclusion assertion.
\end{proof}

\subsection{A truncated Hankel extension lemma}
Hankel rank is closely related to binary apolarity. The following standard
consequence is included with
its degree bookkeeping, since a possible root at infinity matters here.
\begin{lemma}\label{lem:hankel}
Let $\ell$ be a nonzero linear functional on $\C[X]_{\le2m}$, and
suppose its middle Hankel matrix $(\ell(X^{i+j}))_{0\le i,j\le m}$
has rank $r\le m$.  There is a nonzero polynomial $p$ of degree
$s\le r$ and a functional $\widetilde\ell$ on $\C[X]$, annihilated
by $p$, which agrees with $\ell$ through degree $2m-r+s$.
For $s=0$ this means agreement with the zero functional.
\end{lemma}
\begin{proof}
Write $m_j=\ell(X^j)$ and homogenize the moments to the binary form
\[
 f(U,V)=\sum_{j=0}^{2m}\binom{2m}{j}m_jU^{2m-j}V^j.
\]
Its catalecticant maps have Hankel matrices in the divided-power
normalization.  The apolar ideal of a binary form is a complete
intersection with generator degrees $d,2m+2-d$, where $d\le m+1$;
see \cite[Proposition 1.7]{BS}.  Its Hilbert function in degree $m$
is $d$, so $d=r$ when the middle rank is at most $m$.  The case where
$f$ is a pure power has the same conclusion directly, with $r=1$.

A degree-$r$ homogeneous annihilator, dehomogenized, is a polynomial
$p(X)=\sum_{j=0}^sp_jX^j$ with $s\le r$ and $p_s\ne0$.  Its
catalecticant equations are
\[
 \sum_{j=0}^sp_jm_{j+k}=0\qquad(0\le k\le2m-r).
\]
Extend the first $s$ moments by this recurrence.  The resulting
functional is annihilated by $p$, and the displayed equations force
agreement through degree $2m-r+s$.  If $s=0$, those equations say
directly that all moments in the asserted range are zero.  The drop
$r-s$ is exactly what permits homogeneous factors at infinity; no
flat-extension assumption has been inserted. The Lean development proves
the needed short-annihilator and recurrence-extension statements directly;
it does not assume an apolarity theorem as an external axiom.
\end{proof}

\begin{theorem}\label{thm:rank}
Let $m\ge4$ and $2m\le\beta+3$.  For arbitrary complex numbers
$c_3,c_5,c_7,c_9$, put $\ell=C_4+\sum c_jB_j$.  Then
\[
 \rank(\ell(Y^{i+j}))_{0\le i,j\le m}\ge(m+1)-4.
\]
\end{theorem}
\begin{proof}
Assume the contrary, so its rank $r$ is at most $m-4$.
Apply Lemma~\ref{lem:hankel}, and put $e=r-s$.  The extension agrees
through degree $2m-e\ge m+4\ge8$.  It cannot be zero by
Lemma~\ref{lem:independent}, so $s>0$.
Define a functional on $\C[z]$ by
\[
 U(f)=\widetilde\ell\left(\frac{f(z)-f(-z)}{2z}\right),
\]
reading the even quotient as a polynomial in $X=z^2$.  It is odd and
annihilated by $p(z^2)$, and therefore factors through an algebra of
dimension $2s$.  Let $W(q)=U(A(z)q((z+1)^2))$.
By definition, $W(q)=\widetilde\ell(\mathcal Dq)/2$.
For $0\le k\le2s-1$,
\[
 \deg\mathcal D(X^k)\le2s+8\le2m-(r-s),
\]
since $2m-r-s\ge2(m-r)\ge8$.  These moments vanish by agreement
with $\ell$ and Lemma~\ref{lem:pearson}.  Multiplication by
$(z+1)^2$ in $\C[z]/(p(z^2))$ is an endomorphism of a
$2s$-dimensional space.  Cayley--Hamilton extends the vanishing to
every $k\ge0$.  Hence $U$ is a finite Pearson solution.  Its first
nine induced moments agree with $\ell$, contradicting
Proposition~\ref{prop:finite-restrictions}, since the $C_4$ coordinate
of $\ell$ is one.
\end{proof}

\subsection{Selection of a nonzero minor}
Put
\begin{equation}\label{eq:N-K}
 M=100n,\qquad C=\frac{361}{200},\qquad
 N=\floor{361n/2}+1,\qquad K=N-4.
\end{equation}
Then $2N-2\le\beta$ and $(N-1)/M\le C$.  We leave the final
few permissible polynomial degrees unused to retain the latter exact
inequality.  Define
\begin{equation}\label{eq:E}
 E_0(Y)=1,\qquad
 E_i(Y)=\frac{2\prod_{h=0}^{i-1}(Y-h^2)}{(2i)!}\quad(i\ge1).
\end{equation}
These form a basis in every initial degree range.  Moreover
\begin{equation}\label{eq:E-integer}
 E_i(x^2)=\binom{x+i}{2i}+\binom{x+i-1}{2i}\quad(i\ge1),
\end{equation}
an identity of polynomials, so $E_i(x^2)$ is integer-valued on $\Z$.
One can verify~\eqref{eq:E-integer} by adding the two factorial products:
their common factors give $2x^2\prod_{h=1}^{i-1}(x^2-h^2)/(2i)!$.
Generalized binomial coefficients at negative integer arguments are
integers as well.

Apply Theorem~\ref{thm:rank} to~\eqref{eq:zeta-form}.  The matrix
$(F_4(E_iE_j))_{0\le i,j<N}$ has rank at least $K$.  Hence there
exist sets $I,J\subset\{0,\ldots,N-1\}$ of cardinality $K$ for which
\begin{equation}\label{eq:minor}
 D_n=\det\bigl(F_4(E_iE_j)\bigr)_{i\in I,j\in J}\ne0.
\end{equation}
Under simultaneous rationality it belongs to $\Q$.  The rest of the
argument bounds all these minors uniformly; it imposes no condition
on the chosen $I,J$.

\section{The real estimate on a tilted contour}\label{sec:real}
Contour representations of rational-function linear forms are part of the
classical approximation framework; see \cite{Nesterenko03}. Here only an
upper bound is needed: nonvanishing was proved algebraically in
Section~\ref{sec:rank}, so no saddle-point asymptotic is used to establish it.
\subsection{An integral representation}
Set $a=b+\tfrac12$ and $\kappa=3/5$.  Let $\Gamma$ consist of the
two rays $a+(\kappa\pm i)t$, $t\ge0$, oriented as the left boundary
of the region containing the real integers greater than $b$.  Define
\[
 K_5(z)=\frac1{24}\frac{\mathrm d^4}{\mathrm dz^4}(\pi\cot\pi z).
\]
Its principal part at an integer $m$ is exactly $(z-m)^{-5}$.
\begin{lemma}\label{lem:contour}
For $\deg P\le\beta+3$, up to an orientation sign independent of $P$,
\begin{equation}\label{eq:contour}
 F_4(P)=\frac1{2\pi i}\int_\Gamma w_n(z)P(z^2)K_5(z)\dd z.
\end{equation}
On the two rays, $|K_5(a+(\kappa\pm i)t)|\le c_\kappa e^{-2\pi t}$.
\end{lemma}
\begin{proof}
Close the truncated rays at imaginary heights $\pm T$ by horizontal
segments and a far vertical segment at a half-integer real part
$R_T\asymp(\kappa+1)T$.  The rational factor has no pole in the
enclosed region, and the kernel poles there give precisely the Taylor
coefficients in~\eqref{eq:F}.  The horizontal segments have exponentially
small integrals.  On the far vertical segment the kernel is bounded
near the real axis and decays exponentially away from it; the rational
factor tends to zero at least as $R_T^{-3}$.  The closing integrals
therefore vanish for each fixed $n$ and $P$.

For $t$ bounded away from zero, differentiation of the elementary
exponential formula for cotangent gives the stated kernel estimate.
For small $t$, the ray starts at a half-integer and stays away from
the pole set.  Compactness supplies a bound on that remaining segment.
The conjugate ray has the same modulus.  This also proves absolute
convergence of the improper integrals.
\end{proof}

\subsection{The limiting root-product potential}
Write $\alpha_i=a_i/M$, so the fixed list is
\[
 \alpha=(.67,.67,.67,.67,.67,.63,.60,.58,.57,.54,.52,.49,.47,.44,.42).
\]
Let $\beta_0=\sum\alpha_i-5=361/100=2C$, and set
\[
 \zeta_y=1+(\kappa+i)y,\qquad Z_y=\zeta_y^2,\qquad c_y=1+\kappa y.
\]
The notation $\zeta_y$ in this section denotes a point on a ray, not
the Riemann zeta function.  Define
\begin{align}
 f(u)&=u\log u-u\quad(u>0),\qquad f(0)=0,\notag\\
 J(u,y)&=\tfrac u2\log(u^2+y^2)-u+y\arctan(u/y),\notag\\
 S&=\sum_{i=6}^{15}f(2\alpha_i)-2\sum_{i=1}^{5}f(1-\alpha_i),\label{eq:S}\\
 G(y)&=S+5\{J(c_y+1,y)-J(c_y-1,y)\}\notag\\
 &\hspace{18mm}-\sum_{i=1}^{15}\{J(c_y+\alpha_i,y)-J(c_y-\alpha_i,y)\}.
 \label{eq:G}
\end{align}
Continuous limiting values are used when $y=0$ or $u=0$.
Since $\partial J/\partial u=\log|u+iy|$, one can equivalently write
\begin{equation}\label{eq:G-integral}
 G(y)=S+5\int_{-1}^{1}\log|\zeta_y+t|\dd t
                -\sum_i\int_{-\alpha_i}^{\alpha_i}\log|\zeta_y+t|\dd t.
\end{equation}

\begin{lemma}\label{lem:weight-potential}
Uniformly for $y\ge0$ and $n\ge1$,
\begin{equation}\label{eq:weight-potential}
 \log|w_n(a+M(\kappa+i)y)|
 \le MG(y)+O(\log(M+1)+\log(1+y)).
\end{equation}
\end{lemma}
\begin{proof}
Apply Stirling's formula with logarithmic error (cf.~\cite[\S5.11]{DLMF}) to the finitely many
factorials in $\mathcal N_n$.  The coefficient of $M\log M$ is
$2\sum_{i=6}^{15}\alpha_i-2\sum_{i=1}^{5}(1-\alpha_i)=2\beta_0$.
The root products have the opposite coefficient $-2\beta_0$.
Their remaining Riemann sums give the integrals in~\eqref{eq:G-integral}.
For the denominator intervals, $\Re\zeta_y-\alpha_i\ge1-.67>0$;
monotone integral comparison bounds every endpoint error by
$O(\log(M+1)+\log(1+y))$.  The numerator interval has a possible
endpoint singularity at $y=0,t=-1$.  Splitting off its first half-cell
and integrating $\log t$ explicitly bounds the additional error by
$O(\log(M+1))$.  The half-unit displacement $a-M=1/2$ is handled by
the same comparisons.  Finally the explicit factor $2x$ and the
one extra factor in each $D_a$ contribute only logarithmic terms.
There are a fixed number of products, which proves the uniform bound.
\end{proof}

\section{A bounded-root basis adapted to the ray}\label{sec:basis}
\subsection{The exact finite witness}
The accompanying file \path{numerical/replay/round7_basis_witness.json} specifies
$L=45$ rational complex centers $\xi_l$, the rational radius
$\delta=1/50$, and forty vectors $w^{(j)}=(w_l^{(j)})$ of
nonnegative rational weights.  Each weight vector sums exactly to one;
its common denominator is $10^8$.  The centers and nonzero weights are
also reproduced in Appendix~\ref{app:witness}.  All their root disks lie
in $|z|\le R$, where
\begin{equation}\label{eq:R}
 R=\frac{14697}{5000}=2.9394.
\end{equation}
It suffices to verify
$|\Re\xi_l|+|\Im\xi_l|+\delta\le R$, so no numerical square root
is needed for this check.

Use knots $s_j=C(2j+1)/80$, $0\le j<40$, together with $0,C$.
At $s_j$, let the masses be $\mu_l(s_j)=s_jw_l^{(j)}$.
At zero they are all zero; at $C$ use $\mu_l(C)=Cw_l^{(39)}$.
Interpolate the masses linearly on every interval between consecutive
knots.  Thus $\mu_l(s)\ge0$ and $\sum_l\mu_l(s)=s$ on $[0,C]$.
For $s>0$ put $w_l(s)=\mu_l(s)/s$.
Define the truncated logarithmic potentials
\begin{equation}\label{eq:Kpot}
 \mathcal K_l(y)=\log\max\{|Z_y-\xi_l|,\delta\}
            +\log\max\{|Z_y-\overline{\xi_l}|,\delta\}.
\end{equation}
Truncation at $\delta$ removes logarithmic singularities while preserving
a polynomial construction realizing the resulting upper bound.

\subsection{Construction of actual polynomials}
For $0\le i<N$, let $s=i/M$.  When $i>0$ set
$m_l=\floor{iw_l(s)/2}$; when $i=0$ set every $m_l=0$.
For each $m_l>0$ take the real monic block
\begin{equation}\label{eq:block}
 \bigl[(Y-M^2\xi_l)^{m_l}-(M^2\delta)^{m_l}\bigr]
 \bigl[(Y-M^2\overline{\xi_l})^{m_l}-(M^2\delta)^{m_l}\bigr].
\end{equation}
Let $d=i-2\sum_lm_l$, so $0\le d<2L$, and multiply these blocks
by $(Y-M^2)^d$.  Denote the resulting monic degree-$i$ polynomial
by $P_i(Y)$, and put $q_i(Y)=P_i(Y)/(2i)!$.
Its roots have modulus at most $RM^2$; the filler root $M^2$ also
satisfies this bound.  The construction is algebraic and explicit.
It makes no assumption about the roots of an unknown orthogonal polynomial.

\begin{lemma}\label{lem:basis-potential}
Uniformly for $0\le i<N$ and $y\ge0$,
\begin{align}
 \log|q_i(M^2Z_y)|^2
 \le{}&4i\log M-2\log((2i)!)\notag\\
 &+M\sum_{l=1}^L\mu_l(i/M)\mathcal K_l(y)+O(1+\log(1+y)).
 \label{eq:basis-potential}
\end{align}
The same estimate, with an additional uniform $O(1)$ error, holds
when $M\zeta_y$ is replaced by $a+M(\kappa+i)y$ before squaring.
\end{lemma}
\begin{proof}
For each block use $|u^m-v^m|\le2\max(|u|,|v|)^m$.
In the squared modulus this contributes $2m_l\mathcal K_l(y)$
after the common powers of $M$ are removed.  Replacing $2m_l$ by
$iw_l(s)$ costs a bounded amount when $\mathcal K_l$ is negative,
because $\mathcal K_l\ge2\log\delta$; when it is positive the
replacement only enlarges the upper bound.  There are at most $L$
blocks and fewer than $2L$ filler roots.  Their constants and the
filler factors cost $O(1+\log(1+y))$.  This proves the first assertion.

For the displacement, write
$Z'_y=(\zeta_y+1/(2M))^2$.  On any fixed bounded $y$ interval,
$|Z'_y-Z_y|=O(M^{-1})$.  The function
$\log\max(|z-\xi|,\delta)$ is globally $\delta^{-1}$-Lipschitz,
so multiplying the change by a total degree $O(M)$ costs $O(1)$.
For large $y$, the denominators in the logarithmic difference grow
as $y^2$, whereas $|Z'_y-Z_y|=O((1+y)/M)$; the same uniform bound
follows.  The finitely many filler roots can be bounded directly.
\end{proof}

Let $V(s)$ be a continuous piecewise-linear function obtained by
interpolating upper bounds at the specified knots, each satisfying
\begin{equation}\label{eq:V}
 G(y)-2\pi y+\sum_l\mu_l(s)\mathcal K_l(y)\le V(s)
 \qquad(y\ge0).
\end{equation}
Such interpolation remains valid at intermediate $s$, because the
masses, rather than the weights themselves, were interpolated linearly.
Section~\ref{sec:real-cert} supplies explicit valid knot bounds.

\begin{proposition}\label{prop:real-rate}
For every $K$-square minor in~\eqref{eq:minor},
\begin{equation}\label{eq:real-rate}
 \limsup_{n\to\infty}M^{-2}\log|D_n|
 \le \int_0^C V(s)\dd s-2C^2\log(2C)+3C^2.
\end{equation}
\end{proposition}
\begin{proof}
Use the positive measure $\dd\mu=|w_n(z)K_5(z)|\,|\dd z|/(2\pi)$
on the two rays.  The preceding estimates and a Laplace upper bound give
\begin{equation}\label{eq:norm-bound}
 \log\int|q_i(z^2)|^2\dd\mu
 \le4i\log M-2\log((2i)!)+MV(i/M)+O(\log(M+1)).
\end{equation}
Uniformity on bounded $y$ follows immediately from the pointwise bounds.
On the unbounded part, the tail estimate in the next section leaves a
factor $e^{-2\pi M(y-5)}$ times a fixed polynomial; integration changes
the logarithm by only $O(\log(M+1))$.

The bilinear moment matrix has entries $\int q_iq_j\dd\nu$ with
$|\dd\nu|=\dd\mu$.  The integral form of Cauchy--Binet, followed by
Cauchy--Schwarz and the Gram determinant inequality, bounds any minor
by the product of the selected row and column $L^2(\mu)$ norms.
Deleting four indices on each side from the sum of the explicit bounds
in~\eqref{eq:norm-bound} changes it by $O(M\log(M+1))$.
Stirling and Riemann summation give
\[
 M^{-2}\sum_{i=0}^{N-1}\{4i\log M-2\log((2i)!)\}
 \longrightarrow\int_0^C\{-4s\log(2s)+4s\}\dd s
 =-2C^2\log(2C)+3C^2.
\]
Thus~\eqref{eq:real-rate} holds first in the $q$ basis.

For completeness, consider the change to the arithmetic basis $E_i$.
After $Y=M^2Z$, the monic versions of both bases have all roots in a
fixed disk, of radius at most $R_0=\max(R,C^2,1)$.
Their coefficient matrices are unitriangular, with entries of degree
drop $h$ bounded by $\binom{i}{h}R_0^h\le(NR_0)^h$.
Inverting a unitriangular matrix by degree-drop compositions gives
inverse entries bounded by $(2NR_0)^h$.  Restoring the powers of $M$
and factorial diagonal factors still bounds every entry of each
change-of-basis matrix and its inverse by $\exp(O(N\log(N+1)))$.

Let $T$ be the change from $q_i$ to $E_i$.  Its diagonal is
$(1,2,\ldots,2)$, so $\det T=2^{N-1}$.  Jacobi's complementary-minor
identity expresses a $(N-4)$-minor of $T$ as $\det T$ times a
four-by-four minor of $T^{-1}$, up to sign.  Each is therefore
$\exp(O(M\log(M+1)))$ in absolute value.  Applying Cauchy--Binet
on both sides of the moment matrix transfers the bound to every
$E$-minor.  There are only $\binom N4^2$ terms, which add a smaller
logarithmic error.  This proves the proposition without any
assumption about principal minors.
\end{proof}

\section{Certification of the real exponent}\label{sec:real-cert}
\subsection{Cell bounds on the finite interval}
The numerical data are checked using Arb ball arithmetic \cite{Arb} at
256-bit precision, together with exact rational checks of directed rounding.  It encloses every
cell; sampled maxima are not used as rigorous suprema.
For $y>0$, differentiation of~\eqref{eq:G} gives
\begin{equation}\label{eq:derivative-lower}
 (G(y)-2\pi y)'\ge-L_0,\qquad
 L_0=\kappa\sum_i\log\frac{1+\alpha_i}{1-\alpha_i}
                    +\frac{19\pi}{2}.
\end{equation}
Indeed, a symmetric integral of width $2\alpha$ has ray derivative
\[
 \frac\kappa2\log\frac{(c_y+\alpha)^2+y^2}{(c_y-\alpha)^2+y^2}
       +\arctan\frac{c_y+\alpha}{y}-\arctan\frac{c_y-\alpha}{y}.
\]
For $\alpha=1$ this is nonnegative.  For $\alpha<1$ it is at most
$\kappa\log((1+\alpha)/(1-\alpha))+\pi/2$.
Adding the derivative of $-2\pi y$ proves~\eqref{eq:derivative-lower}.
Thus, on a cell $[l,h]$ of width $\Delta$,
\[
 G(y)-2\pi y\le G(h)-2\pi h+L_0\Delta.
\]
Also the quadratic map $Z_y$ differs from linear interpolation between
$Z_l,Z_h$ by at most $(1+\kappa^2)\Delta^2/4$.  Therefore
\begin{equation}\label{eq:distance-cell}
 \max(|Z_y-\xi|,\delta)\le
 \max(|Z_l-\xi|,|Z_h-\xi|,\delta)+\frac{1+\kappa^2}{4}\Delta^2.
\end{equation}
The same applies to $\overline\xi$.  Endpoint squared distances are
rational; an integer-square-root ceiling gives an upward rational
bound for each distance, and Arb bounds its logarithm.

We use $50{,}000$ cells of width $\Delta=10^{-4}$ on $[0,5]$.
Each bound for $G-2\pi y$ and for each $\mathcal K_l$ is rounded
upward to a multiple of $10^{-8}$.  The distance ceilings use scale
$10^{10}$.  Nonnegative integer weights allow the weighted potential
to be evaluated exactly.  The current verification uses arbitrary-precision integer and rational
arithmetic for the algebraic expressions, squared-distance inequalities,
cutoffs, log-pair bounds and weighted knot bounds. The two finite families
of logarithm and arctangent enclosures are checked separately, as described
in Section~\ref{sec:verification}.

\subsection{The unbounded tail}
Let $D=(5+\sum_i\alpha_i^3)/3$.  Expanding the symmetric integral
in~\eqref{eq:G-integral} in powers of $t/\zeta_y$ gives, for
$|\zeta_y|>1$,
\begin{equation}\label{eq:G-tail}
 G(y)\le S-2\beta_0\log|\zeta_y|
                         +\frac{D}{|\zeta_y|^2-1}.
\end{equation}
For example, the absolute value of the correction to
$2\alpha\log|\zeta_y|$ is bounded by
$\sum_{h\ge1}\alpha^{2h+1}/(h(2h+1)|\zeta_y|^{2h})
\le\alpha^3/[3(|\zeta_y|^2-1)]$.
The root bound~\eqref{eq:R} gives
\[
 \sum_l\mu_l(s)\mathcal K_l(y)
 \le2s\log(|\zeta_y|^2+R)
 \le4s\log|\zeta_y|+\frac{2sR}{|\zeta_y|^2}.
\]
Since $s\le C=\beta_0/2$ and $|\zeta_y|\ge y$, the full expression
in~\eqref{eq:V}, for $y\ge5$, is at most
\begin{equation}\label{eq:tail-constant}
 S+\frac D{24}+\frac{2CR}{25}-10\pi-2\pi(y-5).
\end{equation}
The constant part has the certified upper bound $-33.82285507$.
It is included in every knot maximum, and is below all active maxima.
The final negative linear term justifies the integrability assertion
used in~\eqref{eq:norm-bound}.

\subsection{Exact output and its use}
The full-grid linkage checks the upper potential bound at all 42 knots
(including $0$ and $C$) on each of the $50{,}000$ cells. The infinite tail
is bounded analytically. Trapezoidal integration of the resulting
piecewise-linear knot majorant is exact. The output is
\begin{align}
 \int_0^C V(s)\dd s
 &\le-\frac{13277632257660676586509793}
                 {1024000000000000000000000},\label{eq:V-integral}\\
 \int_0^CV(s)\dd s-2C^2\log(2C)+3C^2
 &\le-\frac{577853339}{50000000}=A_*.
 \label{eq:A-certified}
\end{align}
The final logarithmic constant is bounded by a rational inequality proved
in Lean. The external tables provide the cell-level elementary-function
values; the formal interpolation, tail and determinant estimates establish
their applicability. Equations~\eqref{eq:A-certified} and~\eqref{eq:real-rate}
prove~\eqref{eq:main-upper}. More precisely, the formal theorem produces a
single error sequence $E_n$, independent of the selected minor, with
$E_n/n^2\to0$, such that
\begin{equation}\label{eq:uniform-upper}
 |D_n|\le\exp\bigl(A_*(100n)^2+E_n\bigr).
\end{equation}

\section{Local arithmetic of the minors}\label{sec:local}
Assume simultaneous rationality and fix the positive common denominator
$Q$ in~\eqref{eq:arithmetic-form}. We normalize $v_p(p)=1$ and allow
$v_p(0)=+\infty$ in lower-bound statements. Arithmetic improvements for
hypergeometric linear forms exploit prime divisibility beyond a naive common
denominator; see \cite{Zudilin04,Viola03}. Here the relevant gain is at the
level of square minors of local moment matrices.

\subsection{A bound valid at every prime}
Put $L_p=\lfloor\log_p(800n)\rfloor$.
\begin{lemma}\label{lem:small-primes}
Every entry in the matrix defining~\eqref{eq:minor} satisfies
\begin{equation}\label{eq:small-entry}
 v_p(F_4(E_iE_j))\ge-40L_p-v_p(Q).
\end{equation}
Consequently $v_p(D_n)\ge-40KL_p-Kv_p(Q)$.
\end{lemma}
\begin{proof}
We keep the Laurent order in the estimate rather than bounding all
coefficients separately. At each pole, substitute $x=a+p^{L_p}u$.
All nonzero root differences under consideration have absolute value at
most $800n$. For each prime power $p^e\le800n$, the centered root counts
and the factorial floors in the normalization give a lower bound $-30$
for the weight contribution. Indeed a centered count of radius $b$ differs
from $2b/p^e$ by at most one, and the linear terms cancel against the
normalization floors; the ten normalization floor differences lose at most
ten and the five numerator and fifteen denominator counts lose at most
twenty. The explicit factor $x$ has nonnegative contribution. Summing over
$e$ gives the scaled Laurent bound $-30L_p$ for $w_n$.

For $i=j+1\ge1$ the exact factorization
\[
 E_i(x^2)=\frac{2xD_j(x)}{(2j+2)!}
\]
shows that each basis factor loses at most $3L_p$. At each prime power,
the centered count for $D_j$ minus the factorial floor is at least $-3$;
the factor $x$ again contributes nonnegatively. The constant basis vector
requires no loss. It follows that the coefficient of $(x-a)^{-k}$ in
$w_n(x)E_i(x^2)E_j(x^2)$ has valuation at least $(-36+k)L_p$.
The harmonic factor of order $k+4$ has valuation at least
$-(k+4)L_p$, since its summation indices are at most $800n$.
The fixed rational zeta coefficients cost at most $v_p(Q)$.
Thus each local summand has valuation at least $-40L_p-v_p(Q)$.
Summation over poles preserves the lower bound. Expanding a $K$-square
determinant gives the last assertion.
\end{proof}
This bound, including higher prime powers, will be used on the whole
omitted interval $p\le n/20$. It replaces the earlier, weaker small-prime
estimate in the initial manuscript.

\subsection{Nonzero residue classes}
Let $p>\sqrt{800n}$.  Eventually such primes do not divide $Q$.
All factorial arguments, root differences, harmonic indices, and
basis factorial arguments are then below $p^2$.  Put
\begin{align}
 \nu_p&=\sum_{t\in T}\floor{2tn/p}-10\floor{33n/p},\label{eq:nu}\\
 c_a(r)&=\floor{(a+r)/p}+\floor{(a-r)/p}+1,\notag\\
 \lambda_r&=\nu_p+5c_b(r)-\sum_{i=1}^{15}c_{a_i}(r),\label{eq:lambda}\\
 c_r&=\sum_{t\in T}c_{tn}(r),\qquad 1\le r<p/2.\label{eq:capacity}
\end{align}
The number $c_a(r)$ counts the integers in $[-a,a]$ congruent to $r$
modulo $p$.  The two $x$-classes $r,-r$ form one class for $Y=x^2$.
Group the partial-fraction expression~\eqref{eq:arithmetic-form}
accordingly, obtaining a local functional and its bilinear moment
matrix $M_r$.

\begin{lemma}\label{lem:local-entry}
The rank of $M_r$ is at most $c_r$.  In the translated monomial basis,
\begin{equation}\label{eq:entry-bad}
 v_p\bigl(M_r((Y-r^2)^i,(Y-r^2)^j)\bigr)
       \ge\lambda_r+i+j-4.
\end{equation}
Let $j_{\max}$ be the largest absolute actual pole in this class.
If $p>a_0+j_{\max}$, the right side improves to
$\lambda_r+i+j+1$.  If in addition there is only one positive
$Y$-node, every entry in an integral polynomial basis has valuation
at least $\nu_p$.
\end{lemma}
\begin{proof}
At each positive absolute pole $j$, the local expression is a linear
combination of derivatives of $P(Y)$ at $Y=j^2$, of orders below $o_j$.
The paired negative pole gives the same $Y$-node.  The sum of these
jet lengths in the class is precisely $c_r$: counting the $r$ class
in the symmetric interval already accounts for both possible signs
of the positive representative.  This proves the rank bound, without
an extra factor of two.

For the entry estimate set $x=r+pu$, or $x=-r+pu$ at the opposite
class.  Factoring out the powers of $p$ from the numerator and
denominator of $w_n$ gives $p^{\lambda_r}$ times a rational function
whose scaled Laurent coefficients at the relevant poles are integral.
Factors not in this residue class are units; distinct scaled poles
have unit differences because the unscaled differences have absolute
value less than $p^2$.  The explicit factor $x$ is a unit as well.
Moreover $x^2-r^2$ has a factor $p$ in these coordinates.
Thus for a polynomial factor $(Y-r^2)^{i+j}$, an order-$k$ Laurent
coefficient in the original $x$ coordinate has valuation at least
$\lambda_r+i+j+k$.  The bracket in~\eqref{eq:arithmetic-form}
loses at most $k+4$ powers of $p$, which proves~\eqref{eq:entry-bad}.

In a good class, namely $p>a_0+j_{\max}$, all harmonic indices are
less than $p$; the brackets are integral.  Every Laurent order is at
least one, giving the improved entry estimate.  Finally, in a good
single-node class, remove the pole factor at that node from the
already cancelled denominator.  All its remaining factors are
$p$-units at the node.  The numerator is an integral polynomial times
$\mathcal N_n$, so all its Laurent coefficients have valuation at
least $v_p(\mathcal N_n)=\nu_p$.  Integral polynomial multiplication
preserves this bound, and so do the good harmonic weights.
\end{proof}

\begin{lemma}\label{lem:local-minor}
Every $m$-minor of $M_r$, in any polynomial family with integral
monomial coefficients, has valuation at least $f_r(m)$, where
\begin{equation}\label{eq:costs}
 f_r(m)=
 \begin{cases}
 m^2+(\lambda_r-5)m,&\text{class not good},\\
 m^2+\lambda_rm,&\text{good, more than one node},\\
 \max\{m^2+\lambda_rm,\nu_pm\},&\text{good, one node}.
 \end{cases}
\end{equation}
The domain is $0\le m\le c_r$; larger minors vanish.
Each displayed integer cost function is convex.
\end{lemma}
\begin{proof}
For an $m$-minor in the translated monomial basis, selected exponents
on either side are distinct nonnegative integers, with sum at least
$m(m-1)/2$.  Summing~\eqref{eq:entry-bad} along a permutation gives
\[
 m(\lambda_r-4)+m(m-1)=m^2+(\lambda_r-5)m.
\]
The good estimate adds $5m$.  The single-node entry bound gives the
additional lower bound $m\nu_p$.  Translation is integral and
unitriangular; Cauchy--Binet transfers each lower bound to every
integral polynomial family.  The rank statement follows from
Lemma~\ref{lem:local-entry}.  Quadratic costs have increasing marginal
increments, and the maximum of the displayed quadratic and affine
costs is convex, including on their integer domains.
\end{proof}

\subsection{A finite dual bound}
The determinant of a sum of local matrices is expanded by assigning its
columns to residue classes and then applying Laplace expansion to the
corresponding groups of rows. Each term is a product of local square minors
of sizes $m_r$, with $\sum_r m_r=K$. For any integer $t$ define
\[
 H_r(t)=\max_{0\le m\le c_r}\{tm-f_r(m)\}.
\]
Then $f_r(m)\ge tm-H_r(t)$ for every admissible $m$, and hence
\begin{equation}\label{eq:dual}
 v_p(D_n)\ge tK-\sum_r H_r(t)-F_p,\qquad
 F_p=\sum_{i\in I}v_p((2i)!)+\sum_{j\in J}v_p((2j)!).
\end{equation}
The factorial cost is applied once after combining the classes. For odd
$p$ the factor two in~\eqref{eq:E} is a unit. This argument needs only a
feasible dual bound, not an optimizer or a proof of optimality.
The class $r=0$ and classes on partition boundaries are included below by
an explicit coarse penalty. Thus no residue class is silently discarded.

The local estimates are initially made for integral numerator polynomials.
The change from translated monomials to those polynomials is integral;
Cauchy--Binet preserves the lower bound. The subsequent division by the
selected row and column factorials gives exactly $F_p$ in~\eqref{eq:dual}.
All of this takes place in the physical variable $Y=x^2$, the same variable
as in the analytic determinant.

\section{Finite arithmetic cells and the prime sum}\label{sec:prime}
\subsection{Uniform rational boxes}
Use the scale
\[
 u=p/n,\qquad v=r/n,\qquad \widehat C=361/2,
\]
so the certified prime range is $1/20<u\le361$. Let
$\mathcal A=(100,67,63,60,58,57,54,52,49,47,44,42)$.
For a rational cell $(l,h]$, partition $[0,h/2]$ into finitely many
rational boxes $[a_B,b_B]$. The actual residue range $[0,u/2]$ is contained
in this interval. For each endpoint $A\in\mathcal A$, the box stores integer
bounds $L_A^\pm,U_A^\pm$ satisfying, at all four corners,
\begin{equation}\label{eq:corner}
 L_A^\pm u\le A\pm v\le(U_A^\pm+1)u,
 \qquad (u,v)\in\{l,h\}\times\{a_B,b_B\}.
\end{equation}
Each difference in these inequalities is affine in $(u,v)$, so the bounds
hold throughout the rectangle. For $a_B<v<b_B$, the nonzero coefficient
of $v$ makes the upper inequality strict. Thus
\[
 L_A^\pm\le\left\lfloor\frac{A\pm v}{u}\right\rfloor\le U_A^\pm
 \quad(l\le u\le h,\ a_B<v<b_B).
\]
This is the corner-to-interval step; a midpoint evaluation would not suffice.
Write $c_A^-=L_A^++L_A^-+1$ and $c_A^+=U_A^++U_A^-+1$.
Normalization floors give an integer $\nu^-\le\nu_p$ valid on the cell.
The box stores integers $c_B,\lambda_B,\nu_B$ such that
\begin{align*}
 c_B&\ge\sum_{A\in T}c_A^+,&
 \lambda_B&\le\nu^-+5c_{100}^- -5c_{67}^+-\sum_{A\in T}c_A^+,
 &\nu_B&\le\nu^-.
\end{align*}
A good-class flag is permitted only when the corresponding pole-location
inequalities hold at every corner. They imply that all harmonic indices
are less than $p$ for an interior residue. A single-node flag additionally
requires $U_{63}^++U_{63}^-\le0$. For a nonempty class this leaves only one
node. These sufficient tests may sacrifice sharpness but never presume
that a classification at one sample point persists across the cell.

Let $g_B(m)$ be the cost~\eqref{eq:costs} formed with the box's lower
parameters and its certified flags. Choose a cell price $t\in\Z$ with
$|t|\le20000$ and nonnegative integers $H_B$ satisfying
\begin{equation}\label{eq:box-dual}
 tm-H_B\le g_B(m)\qquad(0\le m\le c_B,\ m\in\Z).
\end{equation}
A quadratic cost $m^2+am$ reduces this check to the endpoints and integers
nearest $(t-a)/2$. For a cost $\max(m^2+am,\nu m)$, split at
$m=\nu-a$ and check the kink as well. On each resulting interval the
function $tm-g_B(m)$ is either affine or a concave quadratic. This proves
that the finite candidate checks used by the verifier control every
integer $m$ in the stated domain.

\subsection{Boundary residues and the factorial staircase}
There are at most one more boundary points than boxes. For $u\ge1/20$,
every centered count needed here is at most $4002$, every residue capacity
is at most $40020$, and the local weight is at least $-100000$.
Consequently the cost is bounded below by $m^2-100005m$, and for
$|t|\le20000$, $0\le m\le40020$,
\[
 tm-10^{12}\le m^2-100005m.
\]
The fixed penalty $10^{12}$ therefore covers each boundary class, including
$r=0$. This deliberately generous constant affects only an error of lower
asymptotic order.

The number of residues in an interior box is at most
$n(b_B-a_B)+1$. Since $H_B\ge0$, summing the box penalties gives
\[
 \sum_r H_r(t)\le n\sum_B(b_B-a_B)H_B+R_c,\qquad
 R_c=10^{12}(J_c+1)+\sum_BH_B,
\]
where $J_c$ is the number of boxes in the cell. Put $k=\lfloor361/l\rfloor$.
For $u>l$, and $800n<p^2$, the full factorial cost satisfies
\begin{equation}\label{eq:factorial-cost}
 F_p\le n\left(361k-\frac{lk(k+1)}2\right)+2k.
\end{equation}
To see this, enlarge the selected index sets to all $0\le i<N$ and
write $\lfloor2i/p\rfloor$ as the number of positive integers $j$ with
$jp\le2i$. There are at most $k$ levels; the number of indices at level
$j$ is at most $(361n-jp)/2+1$. Sum these counts on both sides of the
minor and use $p>ln$.

If the cell lower value $L_c\in\Q$ satisfies
\begin{equation}\label{eq:cell-lower}
 L_c\le t\widehat C-\sum_B(b_B-a_B)H_B
                   -361k+\frac{lk(k+1)}2,
\end{equation}
then $|K-\widehat Cn|\le4$, together with~\eqref{eq:dual}, gives
\begin{equation}\label{eq:cell-valuation}
 v_p(D_n)\ge nL_c-E_c,\qquad E_c=R_c+4|t|+2k.
\end{equation}
This holds whenever $l<p/n\le h$, $800n<p^2$, $p\nmid Q$, and the
minor is nonzero. In particular it holds throughout the certified prime
range for $n\ge320000$ and $n\ge20Q$.

\subsection{The exact arithmetic certificate}
The data file \path{numerical/global28/arithmetic_cells.txt} gives
$171{,}504$ adjacent rational cells covering $(1/20,361]$, together with
their integer prices and claimed rational lower values. The verifier
reconstructs $3{,}810{,}331$ boxes, including zero-capacity boxes, and
checks~\eqref{eq:corner}--\eqref{eq:cell-lower} in exact rational arithmetic.
Each cell-integral contribution is rounded downward at scale $10^{-8}$.
The verified sum is bounded below by
\begin{equation}\label{eq:padic-cert}
 \sum_c(h_c-l_c)L_c\ge
 -\frac{1434459186663}{12500000}
 =-114756.73493304>-114800.
\end{equation}
The Lean finite predicate requires only the weaker last lower bound
$-114800$. It also requires the exact cell count, ordered endpoints,
adjacency, finite box inequalities, and price bounds.
Its validity is the single externally checked arithmetic-grid admission.
The transfer from that finite predicate to~\eqref{eq:cell-valuation} is
proved in Lean.

\subsection{Chebyshev functions and omitted primes}
Define
\[
 \theta(x)=\sum_{p\le x}\log p,\qquad
 \psi(x)=\sum_{p^a\le x}\log p.
\]
The prime number theorem gives $\theta(x)\sim x$ and $\psi(x)\sim x$;
the formal proof uses the pinned development \cite{PNT}. For each fixed
cell,
\[
 \frac{\theta(h_cn)-\theta(l_cn)}n\longrightarrow h_c-l_c.
\]
Only finitely many such limits are needed. This avoids a limiting
optimization problem or an exchange of a prime sum with an unknown
continuous minimizer.

For $T\le X$, expanding $\psi$ by prime powers gives
\begin{equation}\label{eq:omitted-primes}
 \sum_{p\le T}\lfloor\log_pX\rfloor\log p
 \le\psi(X)-\theta(X)+\theta(T).
\end{equation}
Indeed the first power contributes $\theta(T)$, and the higher powers
form a sub-sum of $\psi(X)-\theta(X)$. Lemma~\ref{lem:small-primes}
therefore bounds the omitted-prime loss by
\[
 40K\{\psi(800n)-\theta(800n)+\theta(n/20)\}.
\]
After division by $(100n)^2$, this tends to
\begin{equation}\label{eq:small-allowance}
 \frac{40(361/2)}{10000\cdot20}=\frac{361}{10000}=0.0361.
\end{equation}
The higher powers contribute zero to this normalized limit. Above $361n$
all relevant factorial arguments, root differences and harmonic indices
are less than $p$; hence the determinant is $p$-integral apart from the
fixed denominator $Q$. Eventually these primes do not divide $Q$ either.
Over all primes the denominator cost is at most $K\log Q=O(n)$.

\subsection{The eventual global lower bound}
Let $E_G=1+\sum_c|E_c|$, a fixed constant. Combining the preceding
bounds with the rational product formula yields
\begin{align}
 \log|D_n|\ge{}&n\sum_c L_c\{\theta(h_cn)-\theta(l_cn)\}
                   -E_G\theta(361n)\notag\\
 &-40K\{\psi(800n)-\theta(800n)+\theta(n/20)\}
                   -K\log Q.\label{eq:global-log-bound}
\end{align}
The cell-boundary convention is $(l_c,h_c]$, so each prime is counted
exactly once. The fixed error and fixed-denominator terms are $o(n^2)$.
Thus the finite predicate gives
\begin{equation}\label{eq:prime-liminf}
 \liminf_{n\to\infty}\frac{\log|D_n|}{(100n)^2}
 \ge\frac{-114800-361}{10000}=-11.5161.
\end{equation}
The actual checked sum in~\eqref{eq:padic-cert} gives the stronger numerical
value $-1438971686663/125000000000=-11.511773493304$ after the same allowance.
For the final theorem it suffices to use the strict slack in
$-11.5161>-11.52$. We obtain, for every sufficiently large $n$,
\begin{equation}\label{eq:P-certified}
 \log|D_n|\ge P_*(100n)^2,\qquad P_*=-\frac{288}{25}=-11.52.
\end{equation}
The cutoff may depend on $Q$ and the fixed certificate constants; no
numerical value for it is asserted. This is the arithmetic rate used in
the completed proof. The stronger midpoint-allocation rate printed in
the 12 October manuscript is superseded and is not used here.

\section{The product-formula contradiction}\label{sec:contradiction}
For a nonzero rational $q$, unique factorization gives
$q=\pm\prod_pp^{v_p(q)}$, with only finitely many nonzero exponents.
Thus
\begin{equation}\label{eq:product-formula}
 \log|q|=\sum_pv_p(q)\log p.
\end{equation}
This is the arithmetic obstruction to a nonzero rational expression being
too small relative to its denominators and divisibility. The same general
principle underlies rational-approximation proofs such as
\cite{Apery,Beukers}; here it is applied to a selected determinant.

Assume all four target zeta values are rational. The rank argument selects
$D_n\in\Q\setminus\{0\}$, and the corrected-coordinate identification
makes it exactly the minor estimated analytically. For sufficiently large
$n\ge1$, the arithmetic estimate~\eqref{eq:P-certified} and the uniform
analytic estimate~\eqref{eq:uniform-upper} imply
\[
 P_*(100n)^2\le\log|D_n|\le A_*(100n)^2+E_n,
 \qquad E_n/n^2\longrightarrow0.
\]
Taking $n$ large enough that $E_n<n^2$ gives
\[
 \frac{1853339}{5000}n^2=370.6678n^2
 \le E_n<n^2,
\]
which is impossible. This proves Theorem~\ref{thm:main}, using the finite
certificates described above. Neither the argument nor its formalization
identifies which of the four values is irrational, proves irrationality
of an individual value, or establishes a three-value assertion.

The normalized separation is exactly
\[
 P_*-A_* =\frac{1853339}{50000000}=0.03706678.
\]
It compares asymptotic rates normalized by $(100n)^2$; it is not an error
bound for a numerical approximation to a zeta value or a determinant.
The proof needs an eventual cutoff, not its explicit numerical size.

\section{Formal verification and finite computations}\label{sec:verification}
\subsection{The theorem and its explicit hypotheses}
The final formal comparison takes a finite arithmetic grid $G$ with
predicate $G.\mathrm{Valid}$ and an analytic finite-knot certificate as
explicit hypotheses. All transitions from those hypotheses to
irrationality are checked in Lean~4 \cite{Lean4}. The two principal
explicit-hypothesis declarations, in namespace
\path{OddZeta.FourValue}, are
\begin{quote}\small
\path{not_all_rational_of_certificates} and
\path{irrational_one_of_four_of_certificates}.
\end{quote}
They occur in \path{FourValueFinalContradiction.lean}.
Their audited axiom dependencies are contained in
\texttt{propext}, \texttt{Classical.choice} and \texttt{Quot.sound};
neither depends on \texttt{sorryAx}.

The module \path{FourValueFinal.lean} supplies the finite certificates
through the isolated interfaces below. Its final declaration is
\begin{quote}\small
\path{OddZeta.FourValue.irrational_zeta7_or_zeta9_or_zeta11_or_zeta13}.
\end{quote}
It states the disjunction of Mathlib's \texttt{Irrational} predicates for
the real parts of the four classical \texttt{riemannZeta} values.
The proved lemmas \path{realZeta_cast} and \path{realZeta_eq_re} identify
these with $\sum_{m=1}^\infty m^{-k}$ for $k>1$.
Thus the conclusion concerns the usual real zeta values. The supplied
wrapper depends on \texttt{sorryAx}, precisely because the numerical
interfaces have not been replaced by Lean proof terms.

\subsection{The eight admitted finite families}
There are eight marked proof admissions in three project files, not eight
unproved asymptotic arguments. Each family is finite and has explicit data
and an independently runnable verifier.
\begin{center}\small
\begin{tabular}{p{.43\textwidth}cp{.40\textwidth}}
\toprule Interface file & Number & Finite content\\\midrule
\path{FourValueExternalTranscendental.lean}&2&Logarithm and arctangent interval enclosures.\\[5pt]
\path{FourValueExternalArithmetic.lean}&5&Rational identities and inequalities connecting the analytic grid to the knot bounds.\\[5pt]
\path{FourValueExternalGlobalArithmetic.lean}&1&Existence of the finite arithmetic grid satisfying its exact rational predicate.\\
\bottomrule
\end{tabular}
\end{center}
The source audit checks the marked locations and pinned statement hashes.
The log/arctangent checker uses outward-rounded Arb balls at 256-bit
precision and verifies containment in the claimed rational intervals
\cite{Arb}. The other two checkers use arbitrary-precision exact integer
and rational arithmetic. No optimizer output or floating-point sampling
is accepted as a substitute for these checks.

The completed reproduction run verified the following quantities.
\begin{center}\small
\begin{tabular}{p{.33\textwidth}p{.57\textwidth}}
\toprule Check & Verified coverage\\\midrule
Transcendental table & $5{,}700{,}022$ logarithm bounds and $1{,}200{,}000$ arctangent bounds.\\[4pt]
Analytic full grid & $50{,}000$ cells and $2{,}100{,}000$ knot checks, together with the exact argument, distance, cutoff and log-pair inequalities.\\[4pt]
Arithmetic grid & $171{,}504$ cells and $3{,}810{,}331$ reconstructed boxes, including zero-capacity boxes.\\
\bottomrule
\end{tabular}
\end{center}
All checks passed. The software checks establish the finite inputs
externally; they do not remove the admitted dependencies in Lean.
The proofs of contour convergence, rank, selected-minor transfer, valuation
assembly, boundary allowances, prime-sum limits and the final contradiction
are outside these finite admissions.

\subsection{Audit scope and dependencies}
The final axiom audit covers $1{,}873$ named declarations. Among the nine
new declarations in the final integration, only the three supplied-certificate
conclusions depend on \texttt{sorryAx}. The explicit-certificate comparison
has no such dependency. The prime-number-theorem source contains two
unrelated admitted declarations in its Wiener module; the audited
conditional global lower-bound theorem does not depend on them.

The build uses Lean 4.33.1, a pinned Mathlib revision \cite{Mathlib}, and
\texttt{PrimeNumberTheoremAnd} \cite{PNT}. The reproduction script completed
end to end, including fresh numerical checks and a fresh axiom audit.
It reused valid compiled artifacts and is not an empty-cache rebuild of
every historical source. An interrupted earlier run is recorded separately;
its partial log is not treated as successful verification. The final
archive was extracted independently and all 528 manifest entries matched.

The manuscript is an exposition of the construction and its verification
boundary. The exact formal statements, numerical predicates and data are
in the companion Round~29 package \cite{Project}. Formal checking verifies
those statements; editorial correspondence with the exposition and
independent mathematical review remain distinct responsibilities.

\appendix
\section{Exact polynomial-basis witness}\label{app:witness}
The following tables reproduce the finite witness used for the real
bound. Centers are indexed from 1. Every omitted weight is zero,
and every displayed weight numerator is divided by $10^8$. Row $j$
belongs to $s_j=C(2j+1)/80$, with $j$ indexed from 0. The first mass
vector at $s=0$ is zero; the terminal mass vector at $s=C$ uses row 39.
The radius is $1/50$, and the ray slope is $3/5$.
The corresponding data file is
\path{numerical/replay/round7_basis_witness.json} in the proof package.
\small
\begin{longtable}{rcc@{\hspace{12mm}}rcc}
\toprule $l$ & $\Re\xi_l$ & $\Im\xi_l$ & $l$ & $\Re\xi_l$ & $\Im\xi_l$\\ \midrule\endhead
1 & $\frac{103}{100}$ & $0$ & 24 & $\frac{3179}{2500}$ & $\frac{2563}{4000}$ \\[4pt]
2 & $\frac{2649}{2500}$ & $\frac{203}{4000}$ & 25 & $\frac{814}{625}$ & $\frac{177}{250}$ \\[4pt]
3 & $\frac{1287}{1250}$ & $\frac{203}{4000}$ & 26 & $\frac{3231}{2500}$ & $\frac{177}{250}$ \\[4pt]
4 & $\frac{2721}{2500}$ & $\frac{103}{1000}$ & 27 & $\frac{1653}{1250}$ & $\frac{3107}{4000}$ \\[4pt]
5 & $\frac{1323}{1250}$ & $\frac{103}{1000}$ & 28 & $\frac{3281}{2500}$ & $\frac{3107}{4000}$ \\[4pt]
6 & $\frac{2621}{2500}$ & $\frac{103}{1000}$ & 29 & $\frac{1677}{1250}$ & $\frac{847}{1000}$ \\[4pt]
7 & $\frac{679}{625}$ & $\frac{627}{4000}$ & 30 & $\frac{3329}{2500}$ & $\frac{847}{1000}$ \\[4pt]
8 & $\frac{2691}{2500}$ & $\frac{627}{4000}$ & 31 & $\frac{34}{25}$ & $\frac{147}{160}$ \\[4pt]
9 & $\frac{696}{625}$ & $\frac{53}{250}$ & 32 & $\frac{27}{20}$ & $\frac{147}{160}$ \\[4pt]
10 & $\frac{2759}{2500}$ & $\frac{53}{250}$ & 33 & $\frac{861}{625}$ & $\frac{124}{125}$ \\[4pt]
11 & $\frac{57}{50}$ & $\frac{43}{160}$ & 34 & $\frac{3419}{2500}$ & $\frac{124}{125}$ \\[4pt]
12 & $\frac{113}{100}$ & $\frac{43}{160}$ & 35 & $\frac{1743}{1250}$ & $\frac{4267}{4000}$ \\[4pt]
13 & $\frac{1457}{1250}$ & $\frac{327}{1000}$ & 36 & $\frac{3461}{2500}$ & $\frac{4267}{4000}$ \\[4pt]
14 & $\frac{2889}{2500}$ & $\frac{327}{1000}$ & 37 & $\frac{1763}{1250}$ & $\frac{1143}{1000}$ \\[4pt]
15 & $\frac{744}{625}$ & $\frac{1547}{4000}$ & 38 & $\frac{3501}{2500}$ & $\frac{1143}{1000}$ \\[4pt]
16 & $\frac{2951}{2500}$ & $\frac{1547}{4000}$ & 39 & $\frac{891}{625}$ & $\frac{4883}{4000}$ \\[4pt]
17 & $\frac{759}{625}$ & $\frac{56}{125}$ & 40 & $\frac{3539}{2500}$ & $\frac{4883}{4000}$ \\[4pt]
18 & $\frac{3011}{2500}$ & $\frac{56}{125}$ & 41 & $\frac{36}{25}$ & $\frac{13}{10}$ \\[4pt]
19 & $\frac{1547}{1250}$ & $\frac{2043}{4000}$ & 42 & $\frac{143}{100}$ & $\frac{13}{10}$ \\[4pt]
20 & $\frac{3069}{2500}$ & $\frac{2043}{4000}$ & 43 & $\frac{1817}{1250}$ & $\frac{5523}{4000}$ \\[4pt]
21 & $\frac{63}{50}$ & $\frac{23}{40}$ & 44 & $\frac{3609}{2500}$ & $\frac{5523}{4000}$ \\[4pt]
22 & $\frac{5}{4}$ & $\frac{23}{40}$ & 45 & $\frac{3641}{2500}$ & $\frac{1463}{1000}$ \\[4pt]
23 & $\frac{801}{625}$ & $\frac{2563}{4000}$ &  &  &  \\[4pt]
\bottomrule\end{longtable}
\begin{longtable}{r p{0.82\textwidth}}
\toprule Row $j$ & Nonzero pairs $(l,\text{weight numerator})$\\\midrule\endhead
0 & $(2,82721227), (6,17278773)$ \\[2pt]
1 & $(3,41767274), (5,51948099), (8,6284627)$ \\[2pt]
2 & $(3,36687301), (5,38270340), (7,6654700), (8,18387659)$ \\[2pt]
3 & $(2,931974), (3,31825824), (5,32215810), (7,24752571)$ \\[2pt]
 & $(10,10273821)$ \\[2pt]
4 & $(2,2032640), (3,28293195), (5,27867542), (7,22703973)$ \\[2pt]
 & $(9,18251540), (12,851110)$ \\[2pt]
5 & $(1,934855), (3,26741462), (5,25321036), (7,21741305)$ \\[2pt]
 & $(9,16685777), (12,8575565)$ \\[2pt]
6 & $(1,2092402), (3,23879064), (5,23056982), (7,20261790)$ \\[2pt]
 & $(9,16098475), (11,12938511), (14,1672776)$ \\[2pt]
7 & $(1,2866677), (3,21695570), (5,21295468), (7,18985445)$ \\[2pt]
 & $(9,16031508), (11,12340093), (14,6785239)$ \\[2pt]
8 & $(1,3674619), (2,2104956), (3,18108154), (5,19338360)$ \\[2pt]
 & $(7,17834262), (9,15366695), (11,10585185), (12,1514413)$ \\[2pt]
 & $(13,10033804), (16,1439552)$ \\[2pt]
9 & $(1,4251483), (2,3679060), (3,15338833), (5,17760838)$ \\[2pt]
 & $(7,16872551), (9,14802052), (11,12182820), (12,127587)$ \\[2pt]
 & $(13,9802262), (16,5182514)$ \\[2pt]
10 & $(1,4715553), (2,5388454), (3,12689033), (5,16299519)$ \\[2pt]
 & $(7,15994628), (9,14175381), (11,12121266), (13,5220122)$ \\[2pt]
 & $(14,4233932), (15,8118540), (18,1043572)$ \\[2pt]
11 & $(1,4890505), (2,5235437), (3,11800023), (5,15427635)$ \\[2pt]
 & $(7,15256356), (9,13609468), (11,12029010), (13,9162349)$ \\[2pt]
 & $(14,633783), (15,8009533), (18,3945901)$ \\[2pt]
12 & $(1,5239511), (2,7260414), (3,9246049), (5,14149161)$ \\[2pt]
 & $(7,14544455), (9,13145072), (11,11520299), (13,10027622)$ \\[2pt]
 & $(15,3617773), (16,3912547), (17,6713412), (20,623685)$ \\[2pt]
13 & $(1,5303690), (2,7004226), (3,8691797), (5,13501365)$ \\[2pt]
 & $(7,13948432), (9,12691567), (11,11238574), (13,9949950)$ \\[2pt]
 & $(15,6772031), (16,1227567), (17,6688005), (20,2982796)$ \\[2pt]
14 & $(1,5433372), (2,7752766), (3,7413498), (5,12676729)$ \\[2pt]
 & $(7,13393084), (9,12154456), (11,11215013), (13,8696173)$ \\[2pt]
 & $(14,702454), (15,8682393), (17,4193235), (18,1864377)$ \\[2pt]
 & $(19,5465769), (22,356681)$ \\[2pt]
15 & $(1,5441148), (2,7489856), (3,7025015), (5,12165625)$ \\[2pt]
 & $(7,12898415), (9,11774994), (11,10883762), (13,8017436)$ \\[2pt]
 & $(14,1220043), (15,8784428), (17,5768133), (18,543969)$ \\[2pt]
 & $(19,5789136), (22,2198040)$ \\[2pt]
16 & $(1,5509655), (2,8057769), (3,6043590), (5,11505147)$ \\[2pt]
 & $(7,12431750), (9,11475334), (11,10385073), (13,9471491)$ \\[2pt]
 & $(15,6239541), (16,1737798), (17,7502581), (19,4031745)$ \\[2pt]
 & $(20,964013), (21,4509494), (24,135019)$ \\[2pt]
17 & $(1,5486318), (2,7800284), (3,5761395), (5,11089155)$ \\[2pt]
 & $(7,12014643), (9,11129861), (11,10121567), (13,9291683)$ \\[2pt]
 & $(15,6134887), (16,1768982), (17,7559694), (19,4140727)$ \\[2pt]
 & $(20,1061521), (21,5049509), (24,1589774)$ \\[2pt]
18 & $(1,5519445), (2,8283731), (3,4951572), (5,10525182)$ \\[2pt]
 & $(7,11628255), (9,10700682), (11,10168785), (13,8771574)$ \\[2pt]
 & $(15,8337128), (17,5813402), (18,918350), (19,6423532)$ \\[2pt]
 & $(21,4320117), (23,3623664), (26,14581)$ \\[2pt]
19 & $(1,5477711), (2,8034975), (3,4742744), (5,10179873)$ \\[2pt]
 & $(7,11270047), (9,10399197), (11,9919006), (13,8597670)$ \\[2pt]
 & $(15,8222551), (17,5182555), (18,1476559), (19,6632417)$ \\[2pt]
 & $(21,3589416), (22,745034), (23,4396828), (26,1133417)$ \\[2pt]
20 & $(1,5478612), (2,8333468), (3,4157045), (5,9732126)$ \\[2pt]
 & $(7,10931495), (9,10195342), (11,9453527), (13,8714464)$ \\[2pt]
 & $(15,7866451), (17,7033626), (19,6110945), (21,5171034)$ \\[2pt]
 & $(23,4074387), (25,2747478)$ \\[2pt]
21 & $(1,5434827), (2,8217991), (3,3907660), (4,81110)$ \\[2pt]
 & $(5,9346406), (7,10599860), (9,9942337), (11,9141107)$ \\[2pt]
 & $(13,8712794), (15,7554028), (17,7143466), (19,3927838)$ \\[2pt]
 & $(20,1858987), (21,5827491), (23,2840469), (24,825409)$ \\[2pt]
 & $(25,3848818), (28,789402)$ \\[2pt]
22 & $(1,5385122), (2,8071441), (3,3726287), (4,1130925)$ \\[2pt]
 & $(5,8254558), (7,10104381), (9,9685718), (11,8912566)$ \\[2pt]
 & $(13,8583731), (15,7398836), (17,7134328), (19,5158107)$ \\[2pt]
 & $(20,703211), (21,5703957), (23,4149062), (25,3851093)$ \\[2pt]
 & $(27,2046677)$ \\[2pt]
23 & $(1,5359796), (2,8496160), (3,3140625), (4,2421511)$ \\[2pt]
 & $(5,6837689), (7,9570351), (9,9330800), (11,9008882)$ \\[2pt]
 & $(13,7936776), (15,7742518), (17,5246871), (18,1237670)$ \\[2pt]
 & $(19,6612926), (21,4318262), (22,520325), (23,5204132)$ \\[2pt]
 & $(25,2694599), (26,417925), (27,3374361), (30,527821)$ \\[2pt]
24 & $(1,5307703), (2,8282826), (3,3033861), (4,2360943)$ \\[2pt]
 & $(5,6655821), (7,9330328), (9,9112318), (11,8817801)$ \\[2pt]
 & $(13,7790190), (15,7625972), (17,5217375), (18,1205092)$ \\[2pt]
 & $(19,6579411), (21,4436785), (22,451807), (23,5279672)$ \\[2pt]
 & $(25,3363612), (27,3619235), (29,1113416), (30,415832)$ \\[2pt]
25 & $(1,5269950), (2,8429425), (3,2682352), (4,2490952)$ \\[2pt]
 & $(5,6245721), (7,9051447), (9,8983686), (11,8316360)$ \\[2pt]
 & $(13,8107070), (15,7033080), (17,6901357), (19,3626989)$ \\[2pt]
 & $(20,2022492), (21,6066180), (23,4121265), (25,4645702)$ \\[2pt]
 & $(27,2344848), (28,361643), (29,2945574), (32,353907)$ \\[2pt]
26 & $(1,5219043), (2,8307505), (3,2541640), (4,2473054)$ \\[2pt]
 & $(5,6037956), (7,8828739), (9,8783144), (11,8144317)$ \\[2pt]
 & $(13,7963697), (15,6920574), (17,6823504), (19,3533818)$ \\[2pt]
 & $(20,2069549), (21,6076042), (23,4155588), (25,4769970)$ \\[2pt]
 & $(27,2901825), (29,3212110), (32,1237925)$ \\[2pt]
27 & $(1,5178627), (2,8632525), (3,2099499), (4,4084683)$ \\[2pt]
 & $(5,4493029), (7,8277223), (9,8481472), (11,8256160)$ \\[2pt]
 & $(13,7354961), (15,7268515), (17,6306482), (19,6150365)$ \\[2pt]
 & $(21,5214648), (23,4986099), (25,3218230), (26,800472)$ \\[2pt]
 & $(27,3872514), (29,1194911), (30,1362787), (31,2532046)$ \\[2pt]
 & $(34,234752)$ \\[2pt]
28 & $(1,5177812), (2,8610902), (3,1861586), (4,4047782)$ \\[2pt]
 & $(5,4304217), (7,8078381), (9,8300920), (11,8088021)$ \\[2pt]
 & $(13,7229422), (15,7149842), (17,6233399), (19,6089184)$ \\[2pt]
 & $(21,5201055), (23,4991424), (25,3094191), (26,972433)$ \\[2pt]
 & $(27,4008111), (29,1266353), (30,1427948), (31,2889126)$ \\[2pt]
 & $(34,977891)$ \\[2pt]
29 & $(1,5169892), (2,8904105), (3,1438511), (4,4172288)$ \\[2pt]
 & $(5,3912430), (7,7848789), (9,8198859), (11,7694092)$ \\[2pt]
 & $(13,7375751), (15,6778547), (17,6399390), (19,5793159)$ \\[2pt]
 & $(21,5381328), (23,2454609), (24,2177412), (25,4764551)$ \\[2pt]
 & $(27,2928266), (28,397153), (29,3674810), (31,2159711)$ \\[2pt]
 & $(33,2224591), (36,151756)$ \\[2pt]
30 & $(1,5166191), (2,8803676), (3,1275923), (4,4099017)$ \\[2pt]
 & $(5,3796092), (7,7677629), (9,8032620), (11,7548618)$ \\[2pt]
 & $(13,7247870), (15,6674699), (17,6315770), (19,5735518)$ \\[2pt]
 & $(21,5346805), (23,2470862), (24,2156983), (25,4781846)$ \\[2pt]
 & $(27,2702879), (28,663830), (29,3843384), (31,2015417)$ \\[2pt]
 & $(32,220511), (33,2635668), (36,788192)$ \\[2pt]
31 & $(1,5153398), (2,9065568), (3,900404), (4,4211756)$ \\[2pt]
 & $(5,3438298), (7,7473363), (9,7865941), (11,7404077)$ \\[2pt]
 & $(13,7121802), (15,6571712), (17,6232564), (19,5677644)$ \\[2pt]
 & $(21,5319773), (23,4765570), (25,2675817), (26,1615538)$ \\[2pt]
 & $(27,4124016), (29,2280247), (30,834382), (31,3266224)$ \\[2pt]
 & $(33,1954935), (35,1958100), (38,88871)$ \\[2pt]
32 & $(1,5145966), (2,8965917), (3,759115), (4,4140984)$ \\[2pt]
 & $(5,3338389), (7,7318239), (9,7713596), (11,7269505)$ \\[2pt]
 & $(13,7001978), (15,6472380), (17,6150066), (19,5616956)$ \\[2pt]
 & $(21,5277856), (23,4747467), (25,2706152), (26,1590380)$ \\[2pt]
 & $(27,4148693), (29,1822856), (30,1327459), (31,3449187)$ \\[2pt]
 & $(33,1511136), (34,484079), (35,2400301), (38,641343)$ \\[2pt]
33 & $(1,5138747), (2,8872377), (3,623751), (4,4073940)$ \\[2pt]
 & $(5,3242973), (7,7170400), (9,7568129), (11,7140624)$ \\[2pt]
 & $(13,6886781), (15,6376346), (17,6069733), (19,5557189)$ \\[2pt]
 & $(21,5236328), (23,4734501), (25,4390587), (27,1596761)$ \\[2pt]
 & $(28,2184206), (29,3897304), (31,2664882), (32,42518)$ \\[2pt]
 & $(33,2979153), (35,1755833), (37,1749266), (40,47671)$ \\[2pt]
34 & $(1,5127640), (2,8779983), (3,497554), (4,4009124)$ \\[2pt]
 & $(5,3152160), (7,7028339), (9,7427820), (11,7015729)$ \\[2pt]
 & $(13,6774425), (15,6281685), (17,5989296), (19,5495356)$ \\[2pt]
 & $(21,5190082), (23,4708062), (25,4381904), (27,1619456)$ \\[2pt]
 & $(28,2175387), (29,3929451), (31,2380151), (32,372144)$ \\[2pt]
 & $(33,3136480), (35,1169410), (36,628204), (37,2195843)$ \\[2pt]
 & $(40,534315)$ \\[2pt]
35 & $(1,5109553), (2,9012433), (3,175826), (4,4456563)$ \\[2pt]
 & $(5,2593737), (7,6779302), (9,7283678), (11,6887305)$ \\[2pt]
 & $(13,6662450), (15,6187022), (17,5909949), (19,5433317)$ \\[2pt]
 & $(21,5144366), (23,4681947), (25,4383592), (27,3930718)$ \\[2pt]
 & $(29,2896996), (30,686799), (31,3266334), (33,1621603)$ \\[2pt]
 & $(34,1058602), (35,2586769), (37,991898), (38,675850)$ \\[2pt]
 & $(39,1563929), (42,19462)$ \\[2pt]
36 & $(1,5095500), (2,8918804), (3,67857), (4,4372576)$ \\[2pt]
 & $(5,2530746), (7,6654079), (9,7154351), (11,6771492)$ \\[2pt]
 & $(13,6557201), (15,6097200), (17,5832206), (19,5371680)$ \\[2pt]
 & $(21,5095831), (23,4650286), (25,4366564), (27,3932263)$ \\[2pt]
 & $(29,2944546), (30,659467), (31,3302942), (33,1672898)$ \\[2pt]
 & $(34,1073678), (35,2675038), (38,1749011), (39,2000550)$ \\[2pt]
 & $(42,453234)$ \\[2pt]
37 & $(1,5084036), (2,8770317), (4,4282346), (5,2497717)$ \\[2pt]
 & $(7,6536848), (9,7031835), (11,6661120), (13,6456499)$ \\[2pt]
 & $(15,6010912), (17,5757170), (19,5311771), (21,5048190)$ \\[2pt]
 & $(23,4618563), (25,4348815), (27,3934360), (29,3648497)$ \\[2pt]
 & $(31,1950894), (32,1240399), (33,3108591), (35,1298616)$ \\[2pt]
 & $(36,1037949), (37,2461204), (39,1463604), (41,1431257)$ \\[2pt]
 & $(44,8490)$ \\[2pt]
38 & $(1,5072352), (2,8521045), (4,4177583), (5,2513121)$ \\[2pt]
 & $(7,6427488), (9,6914955), (11,6554616), (13,6358543)$ \\[2pt]
 & $(15,5926224), (17,5682709), (19,5251258), (21,4998766)$ \\[2pt]
 & $(23,4583594), (25,4326009), (27,3926964), (29,3655079)$ \\[2pt]
 & $(31,1984387), (32,1231095), (33,3146498), (35,977102)$ \\[2pt]
 & $(36,1404513), (37,2616079), (39,685113), (40,799372)$ \\[2pt]
 & $(41,1859665), (44,405870)$ \\[2pt]
39 & $(1,5061914), (2,8277184), (4,4289237), (5,2375656)$ \\[2pt]
 & $(7,6278800), (9,6799795), (11,6449747), (13,6263201)$ \\[2pt]
 & $(15,5843582), (17,5610314), (19,5191927), (21,4950335)$ \\[2pt]
 & $(23,4548739), (25,4303213), (27,3919443), (29,3666736)$ \\[2pt]
 & $(31,3291645), (33,2127166), (34,865133), (35,2767260)$ \\[2pt]
 & $(37,951166), (38,1267720), (39,2196458), (41,1392256)$ \\[2pt]
 & $(43,1305784), (45,5589)$ \\[2pt]
\bottomrule\end{longtable}
\normalsize


\section{Source map and reproducibility}\label{app:reproduction}
\subsection{Principal formal modules}
Paths in this table are relative to the proof package's \texttt{lean/}
directory. The table identifies entry points rather than every dependency.
\begin{longtable}{p{.45\textwidth}p{.47\textwidth}}
\toprule Module & Role\\\midrule\endhead
\path{FourValueUniformNonvanishing.lean}&Uniform rank and nonzero-minor selection.\\[4pt]
\path{FourValueCorrectedZeta.lean}&Classical zeta identity in the corrected coordinates.\\[4pt]
\path{FourValueCorrectedAssembly.lean}&Entrywise and selected-determinant identification.\\[4pt]
\path{FourValueNumericalBridge.lean}&Analytic bound with one uniform $o(n^2)$ error.\\[4pt]
\path{FourValueFullGrid.lean}&Finite knot certificate supplied from the isolated numerical interfaces.\\[4pt]
\path{FourValueGlobalCosts.lean}&Finite dual valuation inequality for the actual corrected matrix.\\[4pt]
\path{FourValueArithmeticCertificate.lean}&Rational box and cell predicates.\\[4pt]
\path{FourValueGlobalPrimeSum.lean}&Product formula and finite weighted Chebyshev sum.\\[4pt]
\path{FourValueGlobalArithmeticRate.lean}&Limit and eventual arithmetic rate.\\[4pt]
\path{FourValueGlobalArithmeticLower.lean}&Nonzero selected rational minor under rationality.\\[4pt]
\path{FourValueFinalContradiction.lean}&Final comparison and real-series identification.\\[4pt]
\path{FourValueFinal.lean}&Supplied-certificate classical zeta conclusion.\\
\bottomrule
\end{longtable}

\subsection{Reproducing the proof-package checks}
Use the companion \path{Four_Value_Lean_Round29.zip}, not the superseded
Round~8 programs. The source package for this manuscript contains the TeX
files and the final report; the much larger proof package contains the
Lean sources, numerical tables and checkers. From the extracted proof
package root, first check the delivered bytes, then install the pinned
Python dependency and run:
\begin{verbatim}
sha256sum --check SHA256_MANIFEST.txt
python3 -m pip install python-flint==0.9.0
bash numerical/final29/reproduce.sh
\end{verbatim}
Lean must be installed using \path{lean/lean-toolchain}. The script obtains
the dependency cache, reruns all three numerical verifiers, builds the
final dependency chain and runs the axiom audit. It clears old success
records before starting and regenerates the combined numerical summary
from the new reports. A failed command stops the run.

The finite data can also be checked separately:
\begin{verbatim}
python3 numerical/external26/verify_values.py \
  numerical/external26/transcendental_values.txt \
  --report numerical/external26/verification_round29.json
python3 numerical/fullgrid27/verify_grid.py \
  --out numerical/fullgrid27/verification_round29.json
python3 numerical/global28/verify_arithmetic.py
\end{verbatim}
Use normal Python execution, not \texttt{python -O}, which disables
assertions. A packaged-evidence and source audit is available as
\begin{verbatim}
python3 numerical/final29/audit_final.py finish
\end{verbatim}
That last command inspects recorded evidence; it does not itself rerun
Lean or the numerical verifiers.

The pinned revisions are
\begin{quote}\small
Lean: \texttt{leanprover/lean4:v4.33.1}.\\
Mathlib: \path{0df444a360eaa60ab8c11dca51a86af692955474}.\\
PrimeNumberTheoremAnd:
\path{650d31264be65f4cd6e70c45d8b25d86d482a761}.
\end{quote}
The exact dependency manifest, data hashes, logs and admission markers are
included in the proof archive. Replacing the eight finite admissions by
Lean proofs would be a further kernel-certification task, not a missing
step in the conditional mathematical reduction.

\subsection{Reproducing this manuscript}
From the manuscript-source directory, run twice:
\begin{verbatim}
pdflatex -interaction=nonstopmode -halt-on-error \
  Four_Odd_Zeta_Manuscript.tex
\end{verbatim}
The source includes the section files, bibliography and exact witness
tables. Historical claims are supported by the references below;
\cite{Project} is the internal computational record and is not presented
as an independently published research article.

\begin{thebibliography}{99}
\bibitem{Apery} R. Ap\'ery,
\emph{Irrationalit\'e de $\zeta(2)$ et $\zeta(3)$},
Ast\'erisque \textbf{61} (1979), 11--13.
\url{https://www.numdam.org/item/AST_1979__61__11_0/}.

\bibitem{BR} K. Ball and T. Rivoal,
\emph{Irrationalit\'e d'une infinit\'e de valeurs de la fonction z\^eta aux entiers impairs},
Invent. Math. \textbf{146} (2001), no.~1, 193--207.
\url{https://doi.org/10.1007/s002220100168}.

\bibitem{Beukers} F. Beukers,
\emph{A note on the irrationality of $\zeta(2)$ and $\zeta(3)$},
Bull. London Math. Soc. \textbf{11} (1979), no.~3, 268--272.
\url{https://doi.org/10.1112/blms/11.3.268}.

\bibitem{BS} M. C. Brambilla and G. Staglian\`o,
\emph{Algebraic boundaries among typical ranks for real binary forms of arbitrary degree},
Found. Comput. Math. \textbf{21} (2021), 1003--1022.
\url{https://doi.org/10.1007/s10208-020-09474-9}.

\bibitem{Lean4} L. de Moura and S. Ullrich,
\emph{The Lean 4 theorem prover and programming language},
in Automated Deduction---CADE 28, Lecture Notes in Computer Science
\textbf{12699}, Springer, 2021, 625--635.
\url{https://doi.org/10.1007/978-3-030-79876-5_37}.

\bibitem{Fischler04} S. Fischler,
\emph{Irrationalit\'e de valeurs de z\^eta (d'apr\`es Ap\'ery, Rivoal, \ldots)},
S\'eminaire Bourbaki, exp.~910, Ast\'erisque \textbf{294} (2004), 27--62.
\url{https://www.numdam.org/item/SB_2002-2003__45__27_0/}.

\bibitem{FSZ} S. Fischler, J. Sprang and W. Zudilin,
\emph{Many odd zeta values are irrational},
Compos. Math. \textbf{155} (2019), no.~5, 938--952.
\url{https://doi.org/10.1112/S0010437X1900722X}.

\bibitem{Arb} F. Johansson,
\emph{Arb: efficient arbitrary-precision midpoint-radius interval arithmetic},
IEEE Trans. Comput. \textbf{66} (2017), no.~8, 1281--1292.
\url{https://doi.org/10.1109/TC.2017.2690633}.

\bibitem{Lai26} L. Lai,
\emph{A note on the number of irrational odd zeta values II},
J. Th\'eor. Nombres Bordeaux \textbf{38} (2026), no.~2, 469--487.
\url{https://doi.org/10.5802/jtnb.1368}.

\bibitem{LY} L. Lai and P. Yu,
\emph{A note on the number of irrational odd zeta values},
Compos. Math. \textbf{156} (2020), no.~8, 1699--1717.
\url{https://doi.org/10.1112/S0010437X20007307}.

\bibitem{Mathlib} The mathlib Community,
\emph{The Lean mathematical library},
in Proceedings of the 9th ACM SIGPLAN International Conference on Certified
Programs and Proofs (CPP 2020), ACM, 2020, 367--381.
\url{https://doi.org/10.1145/3372885.3373824}.

\bibitem{Nesterenko03} Yu. V. Nesterenko,
\emph{Integral identities and constructions of approximations to zeta-values},
J. Th\'eor. Nombres Bordeaux \textbf{15} (2003), no.~2, 535--550.
\url{https://doi.org/10.5802/jtnb.412}.

\bibitem{DLMF} National Institute of Standards and Technology,
\emph{Digital Library of Mathematical Functions},
\S5.11 (Stirling expansions) and \S25.6 (integer arguments of the zeta function).
\url{https://dlmf.nist.gov/5.11}; \url{https://dlmf.nist.gov/25.6}.
Accessed 15 October 2026.

\bibitem{PNT} PrimeNumberTheoremAnd contributors,
\emph{PrimeNumberTheoremAnd}, Lean source development,
revision \path{650d31264be65f4cd6e70c45d8b25d86d482a761}.
\url{https://github.com/AlexKontorovich/PrimeNumberTheoremAnd}.
Exact revision and dependency audit specified in Appendix~\ref{app:reproduction}.

\bibitem{Rivoal00} T. Rivoal,
\emph{La fonction z\^eta de Riemann prend une infinit\'e de valeurs irrationnelles aux entiers impairs},
C. R. Acad. Sci. Paris S\'er. I Math. \textbf{331} (2000), no.~4, 267--270.
\url{https://doi.org/10.1016/S0764-4442(00)01624-4}.

\bibitem{Rivoal02} T. Rivoal,
\emph{Irrationalit\'e d'au moins un des neuf nombres $\zeta(5),\zeta(7),\ldots,\zeta(21)$},
Acta Arith. \textbf{103} (2002), no.~2, 157--167.
\url{https://doi.org/10.4064/aa103-2-5}.

\bibitem{Viola03} C. Viola,
\emph{Birational transformations and values of the Riemann zeta-function},
J. Th\'eor. Nombres Bordeaux \textbf{15} (2003), no.~2, 561--592.
\url{https://doi.org/10.5802/jtnb.414}.

\bibitem{Zudilin01} W. V. Zudilin,
\emph{One of the numbers $\zeta(5),\zeta(7),\zeta(9),\zeta(11)$ is irrational},
Russian Math. Surveys \textbf{56} (2001), no.~4, 774--776.
\url{https://doi.org/10.1070/RM2001v056n04ABEH000427}.

\bibitem{Zudilin04} W. Zudilin,
\emph{Arithmetic of linear forms involving odd zeta values},
J. Th\'eor. Nombres Bordeaux \textbf{16} (2004), no.~1, 251--291.
\url{https://doi.org/10.5802/jtnb.447}.

\bibitem{Zudilin18} W. Zudilin,
\emph{One of the odd zeta values from $\zeta(5)$ to $\zeta(25)$ is irrational.
By elementary means}, SIGMA \textbf{14} (2018), 028, 8~pp.
\url{https://doi.org/10.3842/SIGMA.2018.028}.

\bibitem{Project} \emph{Four-value zeta proof: final contradiction and integration},
Round~29 research archive and final audit, 15 October 2026.
Companion files \path{Four_Value_Lean_Round29.zip} and
\path{Four_Value_Final_Report_2026-10-15.txt}.
Unpublished computational record; contains the formal source, finite data,
verifiers, logs and SHA-256 manifest.
\end{thebibliography}

\end{document}
